% Experimental copy: both finite-mean theorems use the shell-capacity
% endpoint condition.
\documentclass[11pt]{article}

\usepackage[letterpaper,margin=1.1in]{geometry}
\usepackage{amsmath,amssymb,amsthm,mathtools}
\usepackage{microtype}
\usepackage{enumitem}
\usepackage[colorlinks=true,allcolors=blue]{hyperref}
\hypersetup{
  pdftitle={Fixed Points and Short Cycles of Luce Permutations: Poisson Limits and Power-Law Endpoint Singularities}
}

\newtheorem{theorem}{Theorem}[section]
\newtheorem{proposition}[theorem]{Proposition}
\newtheorem{lemma}[theorem]{Lemma}
\newtheorem{corollary}[theorem]{Corollary}
\theoremstyle{remark}
\newtheorem{remark}[theorem]{Remark}
\theoremstyle{plain}
\newtheorem{assumption}[theorem]{Assumption}

\newcommand{\E}{\mathbb E}
\newcommand{\Pp}{\mathbb P}
\newcommand{\R}{\mathbb R}
\newcommand{\one}{\mathbf 1}
\newcommand{\cF}{\mathcal F}
\newcommand{\cL}{\mathcal L}
\newcommand{\eps}{\varepsilon}
\newcommand{\Poi}{\operatorname{Poi}}
\newcommand{\Luce}{\operatorname{Luce}}
\newcommand{\TV}{\operatorname{TV}}
\newcommand{\Var}{\operatorname{Var}}
\newcommand{\Fix}{\operatorname{Fix}}

\title{Fixed Points and Short Cycles of Luce Permutations:\\
Poisson Limits and Power-Law Endpoint Singularities}
\author{}
\date{}

\begin{document}
\maketitle

\begin{abstract}
The Luce, or Plackett--Luce, model generates a nonuniform random permutation
by sampling without replacement with probabilities proportional to fixed
weights.  Its global asymptotic geometry was determined in
\cite{borga2025permutonlocallimitsluce}, where a Poisson law for the fixed
points was conjectured.  We prove this conjecture under weak multiscale
profile and endpoint conditions.  The fixed points converge to an inhomogeneous Poisson
process, and the counts of cycles of any fixed collection of lengths converge
jointly to independent Poisson random variables.  Power-law endpoint
singularities produce a second regime in which the cycle means diverge.  We
prove the corresponding joint central limit theorems and locate the cycles on
logarithmic distance scales.  When both endpoints contribute, the cycle
counts arising from them are asymptotically independent and their means add.
\end{abstract}

\tableofcontents
\clearpage

\section{Introduction}

Among nonuniform laws on permutations, the Luce model is one of the simplest
to define.  Give each label a positive weight and repeatedly choose from the
remaining labels with probability proportional to weight.  The rule is the
same at every step and is consistent under deletion of the labels already
chosen.  Equivalently, attach independent exponential clocks to the labels
and order their ringing times.  The first description makes Luce's law a
natural model of sequential choice; the second makes it a tractable random
permutation.  It is therefore a natural setting in which to ask how a
deterministic inhomogeneity changes the familiar microscopic statistics of a
uniform permutation.

The asymptotic study of this object was initiated in
\cite{borga2025permutonlocallimitsluce}.  Starting from a macroscopic weight
profile, that work identified the limiting permuton, computed pattern
densities, proved local limits for consecutive relative patterns, and
established a central limit theorem for inversions.  These results give a
detailed picture of the global geometry and of the relative order seen from a
typical location.  The same work then asked about a statistic that neither
description sees: the number of fixed points.  It predicted a Poisson law
whose mean is the integral of the Luce permuton density along the diagonal
\cite[Section~5.2]{borga2025permutonlocallimitsluce}.  That question is the
starting point of this paper.

The obstacle is not identifying the candidate mean; it is reaching the scale
on which a fixed point lives.  A fixed point is the exact equality of a label
and its rank.  It lies on a lattice-scale diagonal of zero permuton area,
whereas a consecutive-pattern limit records relative order rather than exact
alignment.  Permutations with the same permuton limit can consequently have
completely different fixed-point behavior.  A general cycle criterion in
\cite{mukherjee2016fixed} reaches the correct scale by assuming a uniform
finite-dimensional local law, but the known Luce limits do not imply that
hypothesis.  The gap is most pronounced in the final positions, where few
clocks remain and the limiting density may be singular.

Our first result proves the prediction in
\cite{borga2025permutonlocallimitsluce} under minimal convergence assumptions
on the weight profile and a multiscale condition at the final labels.  We prove
the spatial statement: the fixed points themselves converge to an
inhomogeneous Poisson process, with intensity given by the diagonal of the
Luce permuton density.

A fixed point is a cycle of length one, so the natural next question is
whether the Poisson behavior persists for every fixed cycle length.  It does.
The counts of cycles of lengths $1,\ldots,L$ converge jointly to independent
Poisson variables, with means given by cyclic traces of the permuton density.
This extension is not a formal product of one-point estimates: a cycle is a
closed chain of exact assignments, and its vertices remain coupled through
the common ranking.

Certain structured failures of endpoint tightness reveal a second regime.
Very small weights ring late, while very large weights ring early.  If a
power-law singularity aligns one of these effects with the corresponding
endpoint labels, then every logarithmic distance scale contributes a constant
number of cycles.  The means are then of order
$\log n$.  We prove the corresponding joint central limit theorems, identify
the constants, and locate the cycles on logarithmic distance scales.  The
cycle counts near the left and right endpoints are asymptotically independent,
and their contributions add in the total count.  The size of the cycle
intensity therefore separates two regimes: bounded intensities give Poisson
limits, while diverging intensities give the corresponding centered and
normalized Gaussian limits.  The Sukhatme profile considered in
\cite{borga2025permutonlocallimitsluce} already belongs to the second regime
for fixed points.

The proof changes character at exactly the places suggested by this
description.  In the bulk, concentration of the exponential race identifies
the predictable probability of an exact match.  At the last few positions,
the remaining total weight may be too small to concentrate relatively; a
counting argument avoids that denominator altogether.  For longer cycles we
resample a finite path of clocks into one background permutation, preserving
the combinatorial constraint that each rank has one preimage.  Finally, at a
power-law endpoint we study the surviving population at large times or the
arrived population at small times.  Logarithmic distance turns the resulting
ratio kernel into a translation-invariant one, and the cycle constants become
its return densities.

Let $\theta_{n,1},\ldots,\theta_{n,n}>0$.  A Luce permutation is obtained by
drawing the labels sequentially without replacement, with probability
proportional to their weights.  Thus, if $\pi_n(r)$ denotes the label drawn at
step $r$, then
\begin{equation}\label{eq:luce-law}
 \Pp(\pi_n=\pi)
 =\prod_{r=1}^n
   \frac{\theta_{n,\pi(r)}}
   {\sum_{j=r}^n\theta_{n,\pi(j)}}.
\end{equation}
We use $R_{n,i}=\pi_n^{-1}(i)$ for the rank of label $i$.  The fixed-point
count and process are
\begin{equation}\label{eq:fixed-process}
 X_n:=\sum_{k=1}^n I_{n,k},\qquad
 I_{n,k}:=\one\{R_{n,k}=k\},\qquad
 \Xi_n:=\sum_{k=1}^n I_{n,k}\,\delta_{k/n}.
\end{equation}
This notation separates the draw order $\pi_n$ from the rank map $R_n$; the
two are inverse permutations, while their fixed points agree.

The exponential-race representation is fundamental.  Let
$E_{n,i}\sim\operatorname{Exp}(\theta_{n,i})$ be independent.  Ordering these
variables increasingly gives the Luce draw order, and
\begin{equation}\label{eq:rank-representation}
 R_{n,i}=1+\sum_{j\ne i}\one\{E_{n,j}<E_{n,i}\}.
\end{equation}
For $\ell\ge1$, let
\begin{equation}\label{eq:cycle-counts}
 C_{n,\ell}:=\frac1\ell
 \sum_{(i_1,\ldots,i_\ell)\in[n]_{\ne}^{\ell}}
 \prod_{a=1}^{\ell}\one\{R_{n,i_a}=i_{a+1}\},
 \qquad i_{\ell+1}:=i_1.
\end{equation}
Thus $C_{n,1}=X_n$.  The draw order and rank map are inverses, so they have
the same cycle counts.

\subsection{Context and relation to earlier work}

The model originates in the theory of stochastic choice.  The
proportional-choice axiom was introduced in~\cite{luce1959individual}, and the
induced law on complete rankings was developed in
\cite{plackett1975analysis}.  The construction is closely related to the
Bradley--Terry model for paired comparisons~\cite{bradley1952rank} and to the
independent-utility viewpoint of~\cite{thurstone1927law}.  Probabilistically,
it is successive size-biased sampling without replacement; the exponential
race representation is discussed in~\cite{pitmantran2015sizebiased}.  The
same distribution is the stationary ranking of the Tsetlin library and has
also prompted exact enumerative questions~\cite{chatterjee2024enumerative}.

Cycle counts provide a classical microscopic benchmark.  Under the uniform
law, the counts of cycles of lengths $1,2,\ldots,L$ converge jointly to
independent Poisson variables with means $1,1/2,\ldots,1/L$; the Ewens family
belongs to the same logarithmic-combinatorial framework
\cite{arratia2003logarithmic,tavare2021magical}.  A nonuniform law need not
retain this behavior.  In the fixed-$q$ Mallows model, for example, bounded
cycle counts can have linear means and Gaussian fluctuations, while the
global cycle geometry undergoes a separate localization transition
\cite{gladkich2018cycle,he2023cycles}.  Short cycles are therefore a sensitive
test of which features of the uniform permutation survive inhomogeneity.

Permutons give the natural global limit theory for permutation sequences
\cite{hoppen2013limits}, but their topology averages away exact lattice
coincidences.  A complementary theorem in~\cite{mukherjee2016fixed} shows
that a uniform local approximation to all fixed-dimensional assignment
probabilities implies independent Poisson limits for short cycles, with means
given by cyclic traces of the limiting density.  This theorem explains the
form of our constants, but does not subsume our results.  Its local hypothesis
is stronger than permuton convergence and is not provided by the Luce bulk or
consecutive-pattern limits, especially near a singular endpoint.  We instead
derive microscopic control from the exponential race under the minimal
profile hypotheses below, and then determine what replaces the finite-mean
Poisson law when an endpoint contributes on every logarithmic scale.

\subsection{Finite-mean regime}

Because multiplying all weights by the same constant does not change the
law, for the finite-mean results we normalize
\begin{equation}\label{eq:normalization}
 \frac1n\sum_{i=1}^n\theta_{n,i}=1.
\end{equation}
Let $f_n$ be the step profile
\begin{equation}\label{eq:step-profile}
 f_n(x)=\theta_{n,i},\qquad (i-1)/n<x\le i/n.
\end{equation}

\begin{assumption}\label{ass:profile}
There is a measurable $f:(0,1)\to(0,\infty)$ such that
\begin{equation}\label{eq:L1-profile}
 \lVert f_n-f\rVert_{L^1(0,1)}\longrightarrow0.
\end{equation}
In particular, $\int_0^1f=1$.
\end{assumption}

Equivalently, under the normalization~\eqref{eq:normalization}, it is enough
to assume $f_n\to f$ almost everywhere and $\int f=1$: Scheffe's lemma then
gives~\eqref{eq:L1-profile}.  The integral condition rules out weight mass
escaping through a vanishing fraction of labels.

Define, for $t\ge0$,
\begin{equation}\label{eq:HFD}
 H(t):=\int_0^1e^{-t f(u)}\,du,\qquad
 F(t):=1-H(t),\qquad
 D(t):=\int_0^1 f(u)e^{-t f(u)}\,du.
\end{equation}
Since $f>0$ almost everywhere, $F$ is continuous and strictly increasing from
$0$ to $1$.  Write $t_x:=F^{-1}(x)$ for $0<x<1$.  The limiting Luce permuton
has density
\begin{equation}\label{eq:rho}
 \rho(x,y)=
 \frac{f(x)e^{-f(x)t_y}}{D(t_y)},
 \qquad (x,y)\in(0,1)^2.
\end{equation}

For the finite-mean theorems, the endpoint condition is most naturally
stated on logarithmic depth shells.  Put
\begin{equation}\label{eq:terminal-depth-notation}
 k_{n,m}:=n-m+1,
 \qquad 1\le m\le n.
\end{equation}
For every integer $j\ge1$, let
\begin{equation}\label{eq:terminal-shells}
 \mathcal D_{n,j}:=
 \left\{m:j\le\log\frac{n}{m}<j+1\right\}.
\end{equation}
All shell cutoffs $J$ below tend to infinity through integers.
When $\mathcal D_{n,j}$ is nonempty, define
\begin{equation}\label{eq:shell-floor}
 r_{n,j}:=\max\mathcal D_{n,j},
 \qquad
 b_{n,j}:=\min_{m\in\mathcal D_{n,j}}\theta_{n,k_{n,m}}.
\end{equation}

\begin{assumption}[Endpoint shell capacity]
\label{ass:fixed-endpoint}
Under the normalization~\eqref{eq:normalization},
\begin{equation}\label{eq:shell-endpoint-condition}
 \lim_{J\to\infty}\limsup_{n\to\infty}
 \sum_{\substack{j\ge J\\\mathcal D_{n,j}\ne\varnothing}}
 \exp\{-b_{n,j}j\}=0.
\end{equation}
\end{assumption}

Since $j\le\log(n/r_{n,j})<j+1$, the condition is unchanged if $j$ in
\eqref{eq:shell-endpoint-condition} is replaced by $\log(n/r_{n,j})$ or by
$j+1$.  To see this, split the shells according as $b_{n,j}\le1$ or
$b_{n,j}>1$; the change costs only a fixed factor in the first case and a
summable geometric error in the second.

Condition~\eqref{eq:shell-endpoint-condition} allows the terminal rates to
vanish and permits different lower envelopes on different logarithmic depth
scales.  A convenient stronger sufficient condition is the existence of
$\gamma>0$, $\eps_0>0$, and $n_0$ such that
\begin{equation}\label{eq:endpoint-lower}
 \theta_{n,k}\ge\gamma
 \quad\text{whenever }n\ge n_0\text{ and }(1-\eps_0)n\le k\le n.
\end{equation}
Another useful sufficient condition, which permits the terminal rates to
vanish, is that for some $\delta>0$ and $J_0$, uniformly over $1\le m\le n$
and all sufficiently large $n$,
\begin{equation}\label{eq:logarithmic-endpoint-lower}
 \theta_{n,n-m+1}\ge(1+\delta)
 \frac{\log\log(en/m)}{\log(en/m)}
 \quad\text{whenever }\log(en/m)\ge J_0.
\end{equation}
Indeed, for $m\in\mathcal D_{n,j}$ one has
$\log(en/m)\in[j+1,j+2)$.  Since $x\mapsto(\log x)/x$ is eventually
decreasing, for all sufficiently large $j$,
\[
 b_{n,j}j
 \ge(1+\delta)\frac{j}{j+2}\log(j+2)
 \ge(1+\delta/2)\log(j+2).
\]
Thus the shell costs are bounded by $(j+2)^{-1-\delta/2}$, whose tail is
summable, and \eqref{eq:logarithmic-endpoint-lower} implies
Assumption~\ref{ass:fixed-endpoint}.
The shell condition is strictly weaker and is the hypothesis used in both
finite-mean theorems below.

\begin{theorem}\label{thm:main-poisson}
Under Assumptions~\ref{ass:profile} and~\ref{ass:fixed-endpoint},
\begin{equation}\label{eq:point-process-limit}
 \Xi_n\ \Longrightarrow\ \operatorname{PRM}
 \bigl(\rho(x,x)\,dx\bigr)
\end{equation}
as random finite point measures on $[0,1]$.  In particular,
\begin{equation}\label{eq:count-poisson}
 d_{\TV}\bigl(\cL(X_n),\Poi(\lambda)\bigr)\longrightarrow0,
 \qquad
 \lambda:=\int_0^1\rho(x,x)\,dx<\infty.
\end{equation}
\end{theorem}

The conclusion is spatial, not only numerical: it records where the fixed
points occur.  In particular, it implies the predicted Poisson law for their
total number.

For $\ell\ge1$, define
\begin{equation}\label{eq:cycle-intensity}
 \lambda_\ell:=\frac1\ell\int_{[0,1]^\ell}
 \rho(x_1,x_2)\rho(x_2,x_3)\cdots\rho(x_\ell,x_1)
 \,dx_1\cdots dx_\ell.
\end{equation}

\begin{theorem}\label{thm:short-cycles}
Under Assumptions~\ref{ass:profile} and~\ref{ass:fixed-endpoint}, for every
fixed $L$,
\begin{equation}\label{eq:cycle-vector-limit}
 (C_{n,1},\ldots,C_{n,L}) \Longrightarrow\
 (Z_1,\ldots,Z_L),
\end{equation}
where $Z_1,\ldots,Z_L$ are independent and
$Z_\ell\sim\Poi(\lambda_\ell)$.  The convergence in
\eqref{eq:cycle-vector-limit} also holds in total variation on
$\mathbb Z_+^L$.
\end{theorem}

When $f\equiv1$, one has $\rho\equiv1$ and
$\lambda_\ell=1/\ell$, recovering the classical uniform-permutation law.
The factor $1/\ell$ in~\eqref{eq:cycle-intensity} accounts for the $\ell$
possible choices of a root on a directed cycle.  Both theorems allow a merely
measurable limiting profile and highly irregular triangular arrays; no
interior continuity or uniform upper and lower bounds are imposed.

\subsection{Power-law endpoint regime}

The endpoint condition in the Poisson theorem cannot simply be omitted.  A
power zero at the right endpoint produces cycles near the upper-right corner
of the permutation matrix; a sufficiently strong power pole at the left
endpoint produces them near the lower-left corner.  In either case, every
logarithmic distance scale from the corner contributes a constant amount.
Adding those scales gives order $\log n$ cycles.

The central limit theorems below are the diverging-mean counterparts of the
finite-mean Poisson results.  If $\mu_n\to\mu<\infty$, the corresponding
limit is $\Poi(\mu)$.  If instead $\mu_n\to\infty$, then the natural
comparison is
\[
 \frac{\Poi(\mu_n)-\mu_n}{\sqrt{\mu_n}}
 \Longrightarrow\mathcal N(0,1).
\]
Our endpoint theorems prove the corresponding CLTs for the actual cycle
counts, jointly across cycle lengths, endpoint windows, and active corners.
The relevant mean is of order $\log n$ in the power-law regime and, for fixed
points, $\log\log n$ at the critical pole.

\begin{assumption}\label{ass:simple-power-profile}
For some continuous $f:(0,1)\to(0,\infty)$,
\[
 \theta_{n,i}=f((i-1/2)/n).
\]
At the right endpoint we allow either a power zero
\[
 f(1-s)=c_Rs^\beta(1+O(s^\eta)),
 \qquad \beta,c_R,\eta>0,
\]
or a finite positive limit.  At the left endpoint we allow either a
nonintegrable power pole
\[
 f(s)=c_Ls^{-\alpha}(1+O(s^\eta)),
 \qquad \alpha>1,\quad c_L,\eta>0,
\]
or a finite positive limit.  At least one of the two power behaviors is
present.
\end{assumption}

Call an endpoint \emph{active} when the corresponding power behavior is
present.  Thus $f$ is bounded above and below on every interval separated
from the active endpoints, and there are no other singularities.

Near an active corner, transition probabilities depend asymptotically on the
ratio of two endpoint distances.  Taking logarithms turns this ratio into a
difference.  The resulting increment density is, for $\gamma,q>0$,
\begin{equation}\label{eq:generic-power-kernel}
 p_{\gamma,q}(w):=\gamma q e^{\gamma w}e^{-qe^{\gamma w}},
 \qquad
 \mathfrak b_{\gamma,\ell}(q):=\frac1\ell p_{\gamma,q}^{*\ell}(0).
\end{equation}
Here $p^{*\ell}$ denotes $\ell$-fold convolution.  Thus
$p_{\gamma,q}^{*\ell}(0)/\ell$ is the return density of a closed $\ell$-step
path, divided by the number of choices of its root.
For an active right or left endpoint, respectively, put
\begin{align}
 q_R&:=\Gamma(1+1/\beta)^\beta,
 &b_{R,\ell}&:=\mathfrak b_{\beta,\ell}(q_R),\label{eq:simple-q-right}\\
 q_L&:=\Gamma(1-1/\alpha)^{-\alpha},
 &b_{L,\ell}&:=\mathfrak b_{\alpha,\ell}(q_L).\label{eq:simple-q-left}
\end{align}
Set $b_{\sigma,\ell}=0$ when endpoint $\sigma$ is inactive, and put
$B_\ell:=b_{R,\ell}+b_{L,\ell}$.

The coefficient $b_{\sigma,\ell}$ is the mean number of $\ell$-cycles per
unit of logarithmic distance at corner $\sigma$.  In particular,
$b_{R,1}=\beta q_Re^{-q_R}$ and $b_{L,1}=\alpha q_Le^{-q_L}$.

\begin{theorem}
\label{thm:power-clt}
Under Assumption~\ref{ass:simple-power-profile}, for every fixed $L$,
\begin{equation}\label{eq:power-clt}
 \left(
 \frac{C_{n,\ell}-B_\ell\log n}
 {\sqrt{B_\ell\log n}}
 \right)_{\ell=1}^L
 \ \Longrightarrow\ (Z_1,\ldots,Z_L),
\end{equation}
where $Z_1,\ldots,Z_L$ are independent standard normal variables.  Moreover,
$\E C_{n,\ell}\sim B_\ell\log n$ for every fixed $\ell$.
\end{theorem}

The $\log n$ cycles are not concentrated among the first or last few labels.
They are spread uniformly in logarithmic distance from an active endpoint.
To make this precise, root a cycle $\mathfrak c$ at its largest label
$v(\mathfrak c)$ and put
\[
 d_L(\mathfrak c):=v(\mathfrak c),
 \qquad d_R(\mathfrak c):=n+1-v(\mathfrak c).
\]
Thus $u_\sigma(\mathfrak c):=\log d_\sigma(\mathfrak c)/\log n$ records its
location relative to corner $\sigma$ on a scale from $0$ to $1$.  For
$I=(a,b]\subset(0,1)$, let
\begin{equation}\label{eq:spatial-cycle-count}
 C_{n,\ell}^{\sigma}(I)
 :=\#\{\ell\text{-cycles }\mathfrak c:u_\sigma(\mathfrak c)\in I\}.
\end{equation}
For fixed points this is simply the logarithmic location of the point.

\begin{theorem}
\label{thm:power-spatial}
Under the assumptions of Theorem~\ref{thm:power-clt}, fix pairwise disjoint
intervals $I_1,\ldots,I_J$ whose closures lie in $(0,1)$, and fix $L$.
Jointly over active corners $\sigma$, lengths $1\le\ell\le L$, and
$1\le j\le J$,
\begin{equation}\label{eq:power-spatial-clt}
 \frac{C_{n,\ell}^{\sigma}(I_j)
   -b_{\sigma,\ell}|I_j|\log n}
  {\sqrt{b_{\sigma,\ell}|I_j|\log n}}
 \ \Longrightarrow\ \mathcal N(0,1),
\end{equation}
with independent limits.  Moreover, for every such interval $I$,
\begin{equation}\label{eq:power-spatial-mean}
 \E C_{n,\ell}^{\sigma}(I)
 =b_{\sigma,\ell}|I|\log n+o(\log n).
\end{equation}
\end{theorem}

In particular, the expected number of cycles with
$n^a<d_\sigma\le n^b$ is
$b_{\sigma,\ell}(b-a)\log n+o(\log n)$.  Different logarithmic intervals,
cycle lengths, and endpoints fluctuate independently.  When both endpoints
are active, their means add to the constant
$B_\ell=b_{L,\ell}+b_{R,\ell}$ in Theorem~\ref{thm:power-clt}.

This location describes the entire cycle, not only its root.  More precisely,
for every $\delta>0$, the expected number of cycles counted above that contain
a vertex whose logarithmic distance from the same endpoint differs from
$\log d_\sigma(\mathfrak c)$ by more than $\delta\log n$ tends to zero.

At the borderline pole $f(s)\asymp s^{-1}$, only fixed points diverge, and
they do so on the smaller scale $\log\log n$.

\begin{theorem}
\label{thm:critical-pole}
Suppose $\theta_{n,i}=f((i-1/2)/n)$ for a continuous positive $f$ satisfying
\[
 f(s)=c_Ls^{-1}(1+O(s^\eta))
 \quad(s\downarrow0)
\]
for some $c_L,\eta>0$, and suppose $f$ has a finite positive limit at $1$.
Then
\begin{equation}\label{eq:critical-fixed-clt}
 \frac{C_{n,1}-\log\log n}{\sqrt{\log\log n}}
 \Longrightarrow\mathcal N(0,1),
 \qquad
 \E C_{n,1}\sim\log\log n.
\end{equation}
For every fixed $\ell\ge2$, one has $\sup_n\E C_{n,\ell}<\infty$.
\end{theorem}

Theorem~\ref{thm:critical-pole} is the CLT corresponding to a diverging mean,
now with $\mu_n=\log\log n$.  The means of the longer cycles remain bounded
at the critical exponent.

\begin{remark}\label{rmk:power-population}
Theorems~\ref{thm:power-clt} and~\ref{thm:power-spatial} remain valid for the
triangular arrays described by Assumptions~\ref{ass:power-endpoints}
and~\ref{ass:power-population} in Section~\ref{sec:power-law}.  The first
assumption requires power asymptotics only for the rates in fixed
neighborhoods of the two endpoints.  Away from those neighborhoods there is
no continuity or pointwise convergence assumption.  The required information
about the remaining rates is instead expressed through
\[
 H_n(t):=\frac1n\sum_{i=1}^ne^{-t\theta_{n,i}}
 \qquad
 G_n(t):=\frac1n\sum_{i=1}^n(1-e^{-t\theta_{n,i}})
 \qquad\text{and}\qquad
 D_n(t):=\frac1n\sum_{i=1}^n\theta_{n,i}e^{-t\theta_{n,i}}.
\]
Here $H_n(t)$ counts clocks still running at a large time, while $G_n(t)$
counts clocks that have rung by a small time.  Assumption
\ref{ass:power-population} gives uniform regular-variation estimates for
these quantities and for $D_n(t)$.

More precisely, if the right endpoint has local exponent $r_1>0$ and
$H_n(t)\sim A_-t^{-1/\beta_{\rm s}}$, then it contributes at order $\log n$
exactly when $\beta_{\rm s}=r_1$.  Its constants are then
$q_R=c_1A_-^{r_1}$ and
$b_{R,\ell}=\mathfrak b_{r_1,\ell}(q_R)$; if
$\beta_{\rm s}>r_1$, its contribution is zero at this order.  Similarly, if
the left endpoint has exponent $r_0<-1$ and
$G_n(t)\sim A_+t^{1/\alpha_{\rm f}}$, then it contributes exactly when
$\alpha_{\rm f}=-r_0$, with
$q_L=c_0A_+^{r_0}$ and
$b_{L,\ell}=\mathfrak b_{-r_0,\ell}(q_L)$; if
$\alpha_{\rm f}>-r_0$, its contribution is zero at order $\log n$.
Additional slow or fast populations elsewhere in the array are allowed and
enter only through $A_-$ or $A_+$.

The general result is stated immediately after Assumption
\ref{ass:power-population}, and the rest of Section~\ref{sec:power-law}
proves it directly under these weaker hypotheses.  Thus the proof requires
no modification: Assumption~\ref{ass:simple-power-profile} is verified as a
special case in the same section.  The critical-pole result has the analogous
extension under~\eqref{eq:critical-arrival-tail}--\eqref{eq:critical-arrival-weight};
its limiting coefficient is $d=c_0/C_+$ as in
\eqref{eq:critical-pole-coefficient}.
\end{remark}

The profile singled out in
\cite[Section~5.2]{borga2025permutonlocallimitsluce} gives a concrete instance
of the endpoint regime.

\begin{corollary}\label{cor:sukhatme}
Let $\pi_n$ be a Luce permutation with the standard Sukhatme weights
\[
 \theta_{n,i}=n-i+1,
 \qquad 1\le i\le n.
\]
Put $p(w):=e^we^{-e^w}$ and
$b_\ell:=p^{*\ell}(0)/\ell$.  For every fixed $L$,
\begin{equation}\label{eq:sukhatme-cycle-clt}
 \left(
 \frac{C_{n,\ell}-b_\ell\log n}{\sqrt{b_\ell\log n}}
 \right)_{\ell=1}^L
 \Longrightarrow (Z_1,\ldots,Z_L),
\end{equation}
where the coordinates on the right are independent standard normal random
variables, and $\E C_{n,\ell}\sim b_\ell\log n$.  The spatial conclusion of
Theorem~\ref{thm:power-spatial} also holds, with a single active right corner
and coefficients $b_\ell$.  In particular,
\[
 \frac{C_{n,1}-e^{-1}\log n}{\sqrt{e^{-1}\log n}}
 \Longrightarrow\mathcal N(0,1),
 \qquad
 b_2=K_0(2)\approx0.113894,
\]
where $K_0$ is the modified Bessel function of the second kind.
\end{corollary}

\begin{proof}
Rescale the weights by $n^{-1}$, which does not change the Luce law.  The
triangular-array theorem in Section~\ref{sec:power-law} applies with
$\delta_1=0$, $\beta=1$, $A_-=c_1=1$, and hence $q_R=1$.  This gives
\eqref{eq:sukhatme-cycle-clt} and the spatial assertion.  The first constant
is $p(0)=e^{-1}$, while
\[
 \frac12p^{*2}(0)
 =\frac12\int_{\R}e^{-e^w-e^{-w}}\,dw=K_0(2).\qedhere
\]
\end{proof}

Thus the standard Sukhatme model is already in the diverging-mean regime for
fixed points, with mean asymptotic to $e^{-1}\log n$.
Corollary~\ref{cor:sukhatme} proves the corresponding fixed-point CLT and
confirms the behavior suggested by the simulations in
\cite[Section~5.2]{borga2025permutonlocallimitsluce}.  More generally, when a
right power zero and a left nonintegrable power pole are both present, cycles
occur near both endpoints and their contributions add.

\section{Predictable fixed-point probabilities}

We expose the permutation in draw order.  At step $k$, only one event can
create a fixed point: label $k$ must still be available and must be chosen
next.  Its conditional probability is explicit even though the fixed-point
indicators themselves are dependent.

Let $\cF_{n,k}=\sigma(\pi_n(1),\ldots,\pi_n(k))$.  Immediately before draw
$k$, define the remaining weight
\begin{equation}\label{eq:remaining-weight}
 W_{n,k}:=\sum_{i=1}^n\theta_{n,i}\one\{R_{n,i}\ge k\}.
\end{equation}
The conditional chance of a fixed point is
\begin{equation}\label{eq:predictable-p}
 p_{n,k}:=\Pp(I_{n,k}=1\mid\cF_{n,k-1})
 =\frac{\theta_{n,k}\one\{R_{n,k}\ge k\}}{W_{n,k}}.
\end{equation}
Both the numerator indicator and the denominator are
$\cF_{n,k-1}$-measurable.  The random measure
\begin{equation}\label{eq:compensator-measure}
 A_n:=\sum_{k=1}^n p_{n,k}\delta_{k/n}
\end{equation}
is the predictable compensator of $\Xi_n$ in draw order.  The fixed-point
problem is therefore reduced to showing that $A_n$ approaches a deterministic
measure and that its largest atom vanishes.  The following criterion records
this reduction.

\begin{lemma}\label{lem:predictable-poisson}
Let $I_{n,k}\in\{0,1\}$ be adapted, let $x_{n,k}$ be deterministic points in
a compact space $\mathsf X$, and let
$p_{n,k}=\E(I_{n,k}\mid\cF_{n,k-1})$.  Suppose that for a finite deterministic
measure $\nu$ on $\mathsf X$,
\begin{equation}\label{eq:poisson-criterion}
 \sum_kp_{n,k}\delta_{x_{n,k}}\Longrightarrow\nu
 \quad\text{in probability},
 \qquad
 \max_kp_{n,k}\longrightarrow0
 \quad\text{in probability}.
\end{equation}
Then $\sum_k I_{n,k}\delta_{x_{n,k}}$ converges to a Poisson random measure
with intensity $\nu$.
\end{lemma}

\begin{proof}
Fix a nonnegative continuous $g$, abbreviate
$g_{n,k}=g(x_{n,k})$ and $q_{n,k}=1-e^{-g_{n,k}}$.  We first stop the array
before either its total compensator exceeds $K$ or one conditional
probability exceeds $\delta<1$.  Formally, define the predictable indicators
recursively by
\[
 H_{n,k}:=\one\left\{p_{n,k}\le\delta,
 \ \sum_{j<k}H_{n,j}p_{n,j}+p_{n,k}\le K\right\}.
\]
Set
$\widetilde I_{n,k}=H_{n,k}I_{n,k}$ and
$\widetilde p_{n,k}=H_{n,k}p_{n,k}$.  Then
\begin{equation}\label{eq:likelihood-martingale}
 L_{n,m}:=\prod_{k\le m}
 \frac{e^{-g_{n,k}\widetilde I_{n,k}}}
 {1-\widetilde p_{n,k}q_{n,k}}.
\end{equation}
Each factor has conditional mean one, so $(L_{n,m})_m$ is a nonnegative
mean-one martingale.  This likelihood martingale replaces the product
factorization available for independent indicators: here the future
probabilities may depend on the preceding outcomes.

For a single factor $Y_{n,k}$ in~\eqref{eq:likelihood-martingale},
\[
 \E(Y_{n,k}^2\mid\cF_{n,k-1})
 =1+\frac{\widetilde p_{n,k}(1-\widetilde p_{n,k})q_{n,k}^2}
 {(1-\widetilde p_{n,k}q_{n,k})^2}
 \le \exp(C_{g,\delta}\widetilde p_{n,k}).
\]
Iteration therefore gives
\begin{equation}\label{eq:likelihood-L2}
 \sup_m\E L_{n,m}^2
 \le \exp(C_{g,\delta}K),
\end{equation}
so the stopped likelihoods are uniformly integrable.  Meanwhile,
\begin{align*}
 \prod_k(1-\widetilde p_{n,k}q_{n,k})
 &=\exp\left\{-\sum_k\widetilde p_{n,k}q_{n,k}
       +O\left(\max_k\widetilde p_{n,k}
                    \sum_k\widetilde p_{n,k}\right)\right\}\\
 &\longrightarrow
 \exp\left\{-\int(1-e^{-g})\,d\nu\right\}
\end{align*}
in probability whenever no term is deleted.  By~\eqref{eq:poisson-criterion},
choose $K>\nu(\mathsf X)+1$ and any fixed $\delta\in(0,1)$.  The probability
that a term is deleted tends to zero.  Multiplying the last convergence by
the uniformly integrable likelihood in~\eqref{eq:likelihood-martingale}, and
then removing the stopping, gives
\[
 \E\exp\left\{-\sum_k g(x_{n,k})I_{n,k}\right\}
 \longrightarrow
 \exp\left\{-\int(1-e^{-g})\,d\nu\right\},
\]
the Laplace functional of $\operatorname{PRM}(\nu)$.
\end{proof}

\section{The interior compensator}

On every compact subinterval of $[0,1)$, a law of large numbers for the
exponential race controls both the time of the $k$th arrival and the total
rate still present at that time.  This identifies the compensator in the
bulk.  The endpoint, where relative concentration can fail, is postponed to
Section~\ref{sec:endpoint}.

Let
\begin{align}
 \widehat F_n(t)&:=\frac1n\sum_{i=1}^n\one\{E_{n,i}\le t\},
 &F_n(t)&:=\E\widehat F_n(t)
 =\frac1n\sum_{i=1}^n(1-e^{-\theta_{n,i}t}),\label{eq:empirical-F}\\
 \widehat D_n(t)&:=\frac1n\sum_{i=1}^n
 \theta_{n,i}\one\{E_{n,i}\ge t\},
 &D_n(t)&:=\E\widehat D_n(t)
 =\frac1n\sum_{i=1}^n\theta_{n,i}e^{-\theta_{n,i}t}.
 \label{eq:empirical-D}
\end{align}
Let $\tau_{n,k}$ be the $k$th order statistic of the exponential variables.
Then $W_{n,k}=n\widehat D_n(\tau_{n,k})$ and
\begin{equation}\label{eq:survival-indicator}
 \one\{R_{n,k}\ge k\}=\one\{E_{n,k}\ge\tau_{n,k}\}.
\end{equation}

\begin{lemma}\label{lem:race-law}
Under Assumption~\ref{ass:profile}, for every finite $T$,
\begin{equation}\label{eq:uniform-race}
 \sup_{0\le t\le T}|\widehat F_n(t)-F(t)|\longrightarrow0,
 \qquad
 \sup_{0\le t\le T}|\widehat D_n(t)-D(t)|\longrightarrow0
\end{equation}
in probability.  Consequently, for every $\alpha<1$,
\begin{equation}\label{eq:uniform-quantile}
 \max_{1\le k\le\alpha n}|\tau_{n,k}-t_{k/n}|\longrightarrow0
 \quad\text{in probability}
\end{equation}
and
\begin{equation}\label{eq:uniform-denominator}
 \max_{1\le k\le\alpha n}
 \left|\frac{W_{n,k}}n-D(t_{k/n})\right|\longrightarrow0
 \quad\text{in probability}.
\end{equation}
\end{lemma}

\begin{proof}
First, $L^1$ convergence implies
\begin{equation}\label{eq:max-weight}
 \frac1n\max_i\theta_{n,i}\longrightarrow0.
\end{equation}
Indeed, the mass of any interval of length $1/n$ under $f$ tends uniformly to
zero, while $\lVert f_n-f\rVert_1\to0$.

The maps $a\mapsto e^{-ta}$ and $a\mapsto ae^{-ta}$ are uniformly Lipschitz
in $a\ge0$ when $t$ ranges over a compact interval; for the second map one may
use $|(1-ta)e^{-ta}|\le1$.  Hence $F_n\to F$ and $D_n\to D$ uniformly on
$[0,T]$.  At a fixed time, Hoeffding's inequality gives concentration of
$\widehat F_n$, while
\[
 \Var(\widehat D_n(t))
 \le \frac1{n^2}\sum_i\theta_{n,i}^2e^{-\theta_{n,i}t}
 \le \frac{\max_i\theta_{n,i}}n.
\]
A finite grid, monotonicity in $t$, and uniform continuity of $F$ and $D$
upgrade the fixed-time bounds to~\eqref{eq:uniform-race}.

Since $F$ is strictly increasing, uniform convergence of the distribution
functions implies uniform convergence of their quantiles on $[0,\alpha]$,
which gives~\eqref{eq:uniform-quantile}.  Substituting the random quantiles in
the second part of~\eqref{eq:uniform-race} gives
\eqref{eq:uniform-denominator}.
\end{proof}

\begin{proposition}\label{prop:interior-compensator}
Under Assumption~\ref{ass:profile}, for every $\alpha<1$ and every continuous
$g:[0,\alpha]\to\R$,
\begin{equation}\label{eq:weighted-compensator}
 \sum_{k\le\alpha n}g(k/n)p_{n,k}
 \longrightarrow
 \int_0^\alpha g(x)\rho(x,x)\,dx
\end{equation}
in probability.  Moreover,
\begin{equation}\label{eq:interior-max-p}
 \max_{k\le\alpha n}p_{n,k}\longrightarrow0
 \quad\text{in probability}.
\end{equation}
\end{proposition}

\begin{proof}
Put $t_k=t_{k/n}$ and $B_{n,k}=\one\{E_{n,k}\ge t_k\}$.  Since
$d_\alpha:=\min_{0\le t\le t_\alpha}D(t)>0$, Lemma~\ref{lem:race-law} gives
\begin{equation}\label{eq:p-first-approx}
 \sum_{k\le\alpha n}g(k/n)p_{n,k}
 =\frac1n\sum_{k\le\alpha n}
 \frac{g(k/n)\theta_{n,k}\one\{E_{n,k}\ge\tau_{n,k}\}}
 {D(t_k)}+o_{\Pp}(1).
\end{equation}
We next replace the survival indicator by $B_{n,k}$.

Fix $\eta>0$.  The contribution of $k\le\eta n$ is bounded by a constant
times $n^{-1}\sum_{k\le\eta n}\theta_{n,k}$, whose limit is
$\int_0^\eta f$.  For $\eta n<k\le\alpha n$, $t_k$ is bounded away from zero.
On the event in~\eqref{eq:uniform-quantile}, the two indicators can differ
only when $E_{n,k}$ lies in a shrinking interval around $t_k$.  Uniformly on
this range,
\[
 \theta_{n,k}\Pp(|E_{n,k}-t_k|\le\delta)
 \le\int_{t_k-\delta}^{t_k+\delta}
       \theta_{n,k}^2e^{-\theta_{n,k}s}\,ds
 \le C_\eta\delta.
\]
First send $n\to\infty$, then $\delta\downarrow0$, and finally
$\eta\downarrow0$.  This proves that the replacement costs $o_{\Pp}(1)$.

The variables $B_{n,k}$ are independent.  On $k>\eta n$, the variance of the
resulting normalized sum is $O_\eta(n^{-1})$, because
$\sup_{a\ge0}a^2e^{-a t_\eta}<\infty$.  The initial $\eta n$ terms are again
controlled by $n^{-1}\sum_{k\le\eta n}\theta_{n,k}$.  Therefore the sum is
asymptotic in probability to its expectation.  It remains to identify that
expectation.  On $[\eta,\alpha]$, view the sum as the integral of a step
function and put $k_n(x)=\lceil nx\rceil$.  Uniformly there,
$t_{k_n(x)/n}\to t_x$.  The map $a\mapsto ae^{-at}$ is globally
one-Lipschitz, while, for $t\ge t_\eta/2$,
$\sup_{a\ge0}a^2e^{-at}<\infty$.  Hence the uniform convergence of the times,
the continuity of $g$ and $D$, and $\lVert f_n-f\rVert_1\to0$ give
\begin{align*}
 \frac1n\sum_{k\le\alpha n}
 \frac{g(k/n)\theta_{n,k}e^{-\theta_{n,k}t_k}}{D(t_k)}
 &\longrightarrow
 \int_0^\alpha
 \frac{g(x)f(x)e^{-f(x)t_x}}{D(t_x)}\,dx\\
 &=\int_0^\alpha g(x)\rho(x,x)\,dx.
\end{align*}
The discarded initial block is bounded by
$C_\alpha\int_0^\eta f_n(x)\,dx$; it vanishes after first taking
$n\to\infty$ and then $\eta\downarrow0$.  This also handles $x=0$.

For the maximum, on the high-probability event from
Lemma~\ref{lem:race-law},
\[
 \max_{k\le\alpha n}p_{n,k}
 \le \frac{2\max_i\theta_{n,i}}{nd_\alpha}=o(1)
\]
by~\eqref{eq:max-weight}.
\end{proof}

\begin{corollary}\label{cor:interior-poisson}
For every $\alpha<1$,
\begin{equation}\label{eq:interior-poisson}
 \Xi_n|_{[0,\alpha]}\Longrightarrow
 \operatorname{PRM}\bigl(\rho(x,x)\one_{[0,\alpha]}(x)\,dx\bigr).
\end{equation}
\end{corollary}

\begin{proof}
Apply Lemma~\ref{lem:predictable-poisson} and
Proposition~\ref{prop:interior-compensator}.
\end{proof}

\section{The last few positions}\label{sec:endpoint}

It remains to rule out fixed points in the last $o(n)$ positions, uniformly
enough to let $\alpha\uparrow1$.  Directly concentrating the remaining total
weight is fragile: near the end its expectation may be tiny, so an absolute
concentration estimate gives no useful relative error.  We instead condition
on the ringing time of the candidate label and count how many other clocks
survive.  The argument has no random denominator.

For this section only, suppress the first subscript $n$.  Let
\begin{equation}\label{eq:mean-survivors}
 S(t):=\sum_{i=1}^ne^{-\theta_i t}=\E\sum_{i=1}^n\one\{E_i>t\}.
\end{equation}

\begin{lemma}\label{lem:rank-integral}
For every label $k$,
\begin{equation}\label{eq:rank-integral}
 \Pp(R_k=k)=
 \int_0^\infty \theta_k e^{-\theta_k t}
 \Pp\left(\sum_{j\ne k}\one\{E_j>t\}=n-k\right)dt.
\end{equation}
\end{lemma}

\begin{proof}
Condition on $E_k=t$.  The rank is $k$ exactly when $n-k$ of the other
exponential variables are larger than $t$.
\end{proof}

For $m=1,\ldots,n$, write $k_m=n-m+1$ and put
\[
 N(t):=\sum_{j=1}^n\one\{E_j>t\},
 \qquad
 N_{-k}(t):=\sum_{j\ne k}\one\{E_j>t\}.
\]
For $x>0$, also write
\begin{equation}\label{eq:Q-function}
 \mathcal Q(x):=-\log(1-e^{-x}).
\end{equation}

The following block estimate is the basic endpoint-capacity inequality.  It
uses the actual survivor curve $S$, and therefore can be substantially
stronger than a bound obtained only from the mean-one normalization.

\begin{proposition}[Block endpoint capacity]
\label{prop:block-endpoint-capacity}
Let $\mathcal B\subseteq[n]$ be a nonempty set of terminal depths, and put
\[
 r_{\mathcal B}:=\max\mathcal B,
 \qquad
 a_{\mathcal B}:=\min_{m\in\mathcal B}\theta_{k_m}.
\]
For $t_{\mathcal B}>0$, set
\[
 B_{\mathcal B}:=S(t_{\mathcal B})-1.
\]
If $B_{\mathcal B}>r_{\mathcal B}$, then
\begin{equation}\label{eq:block-endpoint-capacity}
 \E\sum_{m\in\mathcal B}I_{k_m}
 \le
 |\mathcal B|
 \exp\left\{-\frac{(B_{\mathcal B}-r_{\mathcal B})^2}
                         {2B_{\mathcal B}}\right\}
 +\mathcal Q(a_{\mathcal B}t_{\mathcal B}).
\end{equation}
\end{proposition}

\begin{proof}
For $0\le t\le t_{\mathcal B}$ and $m\in\mathcal B$,
\[
 \E N_{-k_m}(t)\ge S(t)-1\ge B_{\mathcal B}.
\]
Since $m-1<r_{\mathcal B}<B_{\mathcal B}$, the Chernoff lower-tail bound
gives
\[
 \Pp(N_{-k_m}(t)=m-1)
 \le
 \exp\left\{-\frac{(B_{\mathcal B}-r_{\mathcal B})^2}
                         {2B_{\mathcal B}}\right\}.
\]
Lemma~\ref{lem:rank-integral}, followed by integration of the exponential
density, bounds the total contribution from $t\le t_{\mathcal B}$ by the
first term in~\eqref{eq:block-endpoint-capacity}.

For the remaining times, independence gives
\begin{align*}
 \Pp(N(t)=m-1)
 &\ge \Pp(E_{k_m}\le t,\,N_{-k_m}(t)=m-1)\\
 &= (1-e^{-\theta_{k_m}t})
    \Pp(N_{-k_m}(t)=m-1).
\end{align*}
Consequently,
\[
 \theta_{k_m}e^{-\theta_{k_m}t}
 \Pp(N_{-k_m}(t)=m-1)
 \le
 \frac{\theta_{k_m}}{e^{\theta_{k_m}t}-1}
 \Pp(N(t)=m-1).
\]
For fixed $t>0$, the map $a\mapsto a/(e^{at}-1)$ is decreasing.  Moreover,
the integers $m-1$, $m\in\mathcal B$, are distinct.  Hence the sum of the
last display over $m\in\mathcal B$ is at most
\[
 \frac{a_{\mathcal B}}{e^{a_{\mathcal B}t}-1}
 \sum_{m\in\mathcal B}\Pp(N(t)=m-1)
 \le \frac{a_{\mathcal B}}{e^{a_{\mathcal B}t}-1}.
\]
Integrating from $t_{\mathcal B}$ to infinity gives
\[
 \int_{t_{\mathcal B}}^\infty
 \frac{a_{\mathcal B}}{e^{a_{\mathcal B}t}-1}\,dt
 =-\log(1-e^{-a_{\mathcal B}t_{\mathcal B}}),
\]
which proves the claim.
\end{proof}

There is a complementary individual estimate for a label whose very small
rate is badly aligned with its prescribed terminal depth.  Let
$A_{m-1}^{(-k_m)}$ denote the sum of the $m-1$ smallest rates among
$\{\theta_i:i\ne k_m\}$, with $A_0^{(-k_1)}=0$.

\begin{lemma}[Individual slow-label charge]
\label{lem:slow-label-charge}
For every $m$,
\begin{equation}\label{eq:slow-label-charge}
 \Pp(R_{k_m}=k_m)
 \le
 \frac{\theta_{k_m}}
      {\theta_{k_m}+A_{m-1}^{(-k_m)}}.
\end{equation}
\end{lemma}

\begin{proof}
Immediately before the draw at which $m$ clocks remain, condition on the
surviving set.  If it contains $k_m$, the conditional chance that $k_m$ rings
next is its rate divided by the total surviving rate.  The other $m-1$
survivors have total rate at least $A_{m-1}^{(-k_m)}$.  Since a fixed point at
$k_m$ requires both survival to this stage and selection at the next draw,
taking expectations proves~\eqref{eq:slow-label-charge}.
\end{proof}

For a common block $m=1,\ldots,M$, the key deterministic observation is
\begin{equation}\label{eq:two-candidate}
 \sum_{m=1}^M\one\{N_{-k_m}(t)=m-1\}\le2.
\end{equation}
Indeed, if $N(t)=\sum_j\one\{E_j>t\}=s$, then the displayed event can hold only
for $m=s$ when $E_{k_m}>t$, or for $m=s+1$ when $E_{k_m}\le t$.  Thus, at any
fixed time, at most two terminal labels can have the correct survivor count.

\begin{proposition}\label{prop:uniform-endpoint-bound}
Assume~\eqref{eq:normalization}.  Let $1\le M\le\eps_0n$, and suppose
$\theta_{k_m}\ge\gamma$ for $1\le m\le M$.  Put
\begin{equation}\label{eq:B-definition}
 B:=\sqrt{nM},
\end{equation}
and let $s$ be determined by $S(s)=B$.  If $M\le(B-1)/2$ and $s\ge1/\gamma$,
then for a universal $c>0$,
\begin{equation}\label{eq:tail-explicit}
 \E\sum_{m=1}^M I_{k_m}
 \le M e^{-cB}+2e^{-\gamma s}
 \le M e^{-cB}+2\left(\frac{2B}{n}\right)^{\gamma/2}.
\end{equation}
Consequently,
\begin{equation}\label{eq:tail-epsilon}
 \limsup_{n\to\infty}
 \E\sum_{k>(1-\eps)n}I_{n,k}
 \le 2^{1+\gamma/2}\eps^{\gamma/4}
\end{equation}
for all sufficiently small $\eps>0$.
\end{proposition}

\begin{proof}
For $t\le s$ and every $m\le M$,
\[
 \E N_{-k_m}(t)\ge S(t)-1\ge B-1\ge2M.
\]
A Chernoff lower-tail bound for a sum of independent Bernoulli variables gives
\[
 \Pp(N_{-k_m}(t)=m-1)\le e^{-cB}.
\]
Using Lemma~\ref{lem:rank-integral} and integrating the exponential density,
the total contribution of $t\le s$ is at most $Me^{-cB}$.

For $t\ge s\ge1/\gamma$, the map $a\mapsto ae^{-at}$ is decreasing on
$[\gamma,\infty)$.  Hence, by~\eqref{eq:two-candidate},
\begin{align*}
 &\sum_{m=1}^M\int_s^\infty
 \theta_{k_m}e^{-\theta_{k_m}t}
 \Pp(N_{-k_m}(t)=m-1)\,dt\\
 &\hspace{1in}\le
 \int_s^\infty\gamma e^{-\gamma t}
 \sum_{m=1}^M\Pp(N_{-k_m}(t)=m-1)\,dt
 \le2e^{-\gamma s}.
\end{align*}

It remains to bound $s$.  The normalization implies that at least $n/2$
weights are at most $2$.  Thus
\[
 S(t)\ge\frac n2e^{-2t}.
\]
Since $S(s)=B$, we get $s\ge\frac12\log(n/(2B))$, proving the second
inequality in~\eqref{eq:tail-explicit}.

Finally take $M=\lceil\eps n\rceil$.  Then $B/n\to\sqrt\eps$, the first
term in~\eqref{eq:tail-explicit} vanishes, and the required side conditions
hold when $\eps$ is sufficiently small.  This gives~\eqref{eq:tail-epsilon}.
\end{proof}

The particular survivor buffer used in the preceding argument is convenient
but not intrinsic.  A square-root buffer will be useful for both fixed points
and short cycles.

\begin{lemma}[Diagonal shell buffer]
\label{lem:diagonal-shell-buffer}
Under Assumption~\ref{ass:fixed-endpoint}, the choice $h_j=\sqrt j$ satisfies
\begin{equation}\label{eq:buffered-shell-condition}
 \lim_{J\to\infty}\limsup_{n\to\infty}
 \sum_{\substack{j\ge J\\\mathcal D_{n,j}\ne\varnothing}}
 \mathcal Q\bigl(b_{n,j}(j-h_j)\bigr)=0.
\end{equation}
Moreover, $h_j\uparrow\infty$, $h_j<j/4$ eventually, and $j-h_j$ is
nondecreasing.
\end{lemma}

\begin{proof}
Write $x_{n,j}=b_{n,j}j$ and omit empty shells.  Put
\[
 y_{n,j}:=b_{n,j}(j-h_j)
 =x_{n,j}\left(1-\frac1{\sqrt j}\right).
 \]
If $x_{n,j}\le\sqrt j$, then $y_{n,j}\ge x_{n,j}-1$ and hence
$e^{-y_{n,j}}\le e e^{-x_{n,j}}$.  Otherwise,
$y_{n,j}>\sqrt j-1$.  Assumption~\ref{ass:fixed-endpoint} therefore gives
\[
 \lim_{J\to\infty}\limsup_n\sum_{j\ge J}e^{-y_{n,j}}
 \le
 e\lim_{J\to\infty}\limsup_n\sum_{j\ge J}e^{-x_{n,j}}
 +\lim_{J\to\infty}\sum_{j\ge J}e^{1-\sqrt j}=0.
\]
The same estimate shows that $y_{n,j}$ tends uniformly to infinity in the
iterated tail.  Hence $\mathcal Q(y_{n,j})\le2e^{-y_{n,j}}$ there, proving
\eqref{eq:buffered-shell-condition}.  The remaining assertions are immediate
from $h_j=\sqrt j$.
\end{proof}

\begin{proposition}[Shell endpoint tightness]
\label{prop:endpoint-bound}
Under the normalization~\eqref{eq:normalization} and
Assumption~\ref{ass:fixed-endpoint},
\begin{equation}\label{eq:shell-tail-tightness}
 \lim_{J\to\infty}\limsup_{n\to\infty}
 \E\sum_{m\le ne^{-J}} I_{n,k_{n,m}}=0.
\end{equation}
Equivalently,
\begin{equation}\label{eq:shell-tail-epsilon}
 \lim_{\eps\downarrow0}\limsup_{n\to\infty}
 \E\sum_{k>(1-\eps)n}I_{n,k}=0.
\end{equation}
\end{proposition}

\begin{proof}
Let $h_j$ be supplied by Lemma~\ref{lem:diagonal-shell-buffer}.  For a
nonempty shell $\mathcal D_{n,j}$, apply
Proposition~\ref{prop:block-endpoint-capacity} with
 \[
 \mathcal B=\mathcal D_{n,j},
 \qquad
 t_{\mathcal B}=\log\frac{n}{r_{n,j}}-h_j.
\]
For all sufficiently large $j$ this time is positive.  Jensen's inequality
and~\eqref{eq:normalization} give
\[
 S(t_{\mathcal B})
 \ge ne^{-t_{\mathcal B}}
 =r_{n,j}e^{h_j}.
\]
It follows that
$B_{\mathcal B}\ge r_{n,j}e^{h_j}-1>r_{n,j}$ for all sufficiently large
$j$, so the hypothesis of Proposition~\ref{prop:block-endpoint-capacity}
holds.  Moreover, for a universal $c>0$,
\[
 \frac{(B_{\mathcal B}-r_{n,j})^2}{2B_{\mathcal B}}
 \ge c r_{n,j}e^{h_j}.
\]
Thus the first term in
\eqref{eq:block-endpoint-capacity} is at most
\[
 |\mathcal D_{n,j}|e^{-c r_{n,j}e^{h_j}}
 \le r_{n,j}e^{-c r_{n,j}e^{h_j}}.
\]
The nonempty shells have distinct maxima.  Since $h_j$ is nondecreasing,
\[
 \sum_{j\ge J}r_{n,j}e^{-c r_{n,j}e^{h_j}}
 \le\sum_{r\ge1}r e^{-c r e^{h_J}}
 \longrightarrow0
 \qquad(J\to\infty),
\]
uniformly in $n$.  The sum of the second terms in
\eqref{eq:block-endpoint-capacity} tends to zero by
\eqref{eq:buffered-shell-condition}, because
$t_{\mathcal B}\ge j-h_j$ and $\mathcal Q$ is decreasing.  Since the union
of the shells $j\ge J$
is precisely the set of depths with $m\le ne^{-J}$, this proves
\eqref{eq:shell-tail-tightness}; the equivalent formulation
\eqref{eq:shell-tail-epsilon} follows by monotonicity.
\end{proof}

\begin{corollary}[Exceptional-shell refinement]
\label{cor:exceptional-shell-refinement}
Assume the normalization~\eqref{eq:normalization}.  For each $n$, let
$\mathcal E_n\subseteq[n]$ be a set of exceptional terminal depths and put
\[
 \mathcal G_{n,j}:=\mathcal D_{n,j}\setminus\mathcal E_n.
\]
When $\mathcal G_{n,j}$ is nonempty, define
\[
 r^{\mathrm G}_{n,j}:=\max\mathcal G_{n,j},
 \qquad
 b^{\mathrm G}_{n,j}
 :=\min_{m\in\mathcal G_{n,j}}\theta_{n,k_{n,m}}.
\]
Let $A_{n,m-1}^{(-k_{n,m})}$ be the sum of the $m-1$ smallest rates after
deleting $\theta_{n,k_{n,m}}$.  Suppose
\begin{align}
 \lim_{J\to\infty}\limsup_{n\to\infty}
 \bigg[
 \sum_{\substack{j\ge J\\\mathcal G_{n,j}\ne\varnothing}}
   e^{-b^{\mathrm G}_{n,j}j}
 \notag+
 \sum_{\substack{m\in\mathcal E_n\\m\le ne^{-J}}}
 \frac{\theta_{n,k_{n,m}}}
      {\theta_{n,k_{n,m}}+A_{n,m-1}^{(-k_{n,m})}}
 \bigg]=0.
 \label{eq:exceptional-shell-condition}
\end{align}
Then the endpoint conclusions
\eqref{eq:shell-tail-tightness}--\eqref{eq:shell-tail-epsilon} hold.
Consequently, Theorem~\ref{thm:main-poisson} remains valid if
Assumption~\ref{ass:fixed-endpoint} is replaced by
\eqref{eq:exceptional-shell-condition}.
\end{corollary}

\begin{proof}
Apply the proof of Lemma~\ref{lem:diagonal-shell-buffer} to the regular
blocks $\mathcal G_{n,j}$.  In the proof of
Proposition~\ref{prop:endpoint-bound}, replace $\mathcal D_{n,j}$,
$r_{n,j}$, and $b_{n,j}$ by their regular-block versions and use the cutoff
$\log(n/r^{\mathrm G}_{n,j})-h_j$.  This cutoff is at least $j-h_j$, so the
same decreasing-$\mathcal Q$ bound makes the regular contribution vanish.
Lemma~\ref{lem:slow-label-charge} bounds the exceptional contribution by the
second sum in~\eqref{eq:exceptional-shell-condition}.  Adding the two proves
the endpoint bounds.  These bounds replace each invocation of
Proposition~\ref{prop:endpoint-bound} in the proof of
Theorem~\ref{thm:main-poisson}, which is otherwise unchanged.
\end{proof}

\begin{remark}\label{rmk:minimal-tail}
The uniform terminal support condition~\eqref{eq:endpoint-lower} implies
Assumption~\ref{ass:fixed-endpoint}.  Indeed, for all sufficiently large
$j$, every label in $\mathcal D_{n,j}$ lies in that terminal interval, so
$b_{n,j}\ge\gamma$ and the shell sum is bounded by the geometric tail
$\sum_{j\ge J}e^{-\gamma j}$.  Proposition
\ref{prop:uniform-endpoint-bound} records the corresponding direct estimate.

At the shell level, \eqref{eq:logarithmic-endpoint-lower} can be weakened
further.  For example,
it is enough that, uniformly on all sufficiently late nonempty shells,
\begin{equation}\label{eq:bertrand-shell-envelope}
 b_{n,j}j
 \ge \log j+(1+\delta)\log\log j
\end{equation}
for some $\delta>0$, because the shell costs are then bounded by
$1/(j(\log j)^{1+\delta})$.  In terms of an individual terminal depth with
$q=\log(n/m)$, the corresponding envelope has order
\[
 \theta_{n,n-m+1}
 \ge
 \frac{\log q+(1+\delta)\log\log q}{q}.
\]
The usual iterated-log refinements of the convergent Bertrand series give
still weaker sufficient envelopes, and sparse bad shells need only have
summable costs.  Provided the remaining shell costs already have uniformly
vanishing tails, a single moving anomalous shell is allowed whenever its
product $b_{n,j}j$ tends to infinity.

The minimum $b_{n,j}$ can nevertheless be spoiled by one harmless ultra-slow
label.  Corollary~\ref{cor:exceptional-shell-refinement} removes such labels
from their regular shells and charges them individually.  This exception
mechanism is proved here only for fixed points.  The short-cycle argument
below uses the lower floor on every label in a shell, so no analogue under
the exceptional-shell refinement is claimed.
\end{remark}

\begin{proof}[Proof of Theorem~\ref{thm:main-poisson}]
For fixed $\alpha<1$, Corollary~\ref{cor:interior-poisson} gives the desired
Poisson limit on $[0,\alpha]$.  Proposition~\ref{prop:endpoint-bound} and
Markov's inequality imply
\begin{equation}\label{eq:tail-tightness}
 \lim_{\alpha\uparrow1}\limsup_{n\to\infty}
 \Pp\bigl(\Xi_n((\alpha,1])>0\bigr)=0.
\end{equation}
It remains to verify that the proposed limiting intensity is finite; weak
convergence alone does not imply convergence of means.  Fix
$\alpha<\beta<1$.  Interior count convergence, followed by truncation and
Fatou's lemma, gives
\begin{align*}
 \int_\alpha^\beta\rho(x,x)\,dx
 &=\E\Poi\left(\int_\alpha^\beta\rho(x,x)\,dx\right)\\
 &\le\liminf_{n\to\infty}\E\Xi_n((\alpha,\beta])
 \le\limsup_{n\to\infty}\E\Xi_n((\alpha,1]).
\end{align*}
Letting $\beta\uparrow1$ and using monotone convergence gives
\[
 \int_\alpha^1\rho(x,x)\,dx
 \le\limsup_{n\to\infty}\E\Xi_n((\alpha,1]).
\]
Proposition~\ref{prop:endpoint-bound} shows that the right side tends to zero
as $\alpha\uparrow1$, first along $\alpha=1-e^{-J}$ and then for arbitrary
$\alpha$ by monotonicity.  Together with finiteness of the intensity on every
compact subinterval of $[0,1)$, this proves both
\begin{equation}\label{eq:intensity-tail-bound}
 \lambda:=\int_0^1\rho(x,x)\,dx<\infty,
 \qquad
 \int_\alpha^1\rho(x,x)\,dx\longrightarrow0.
\end{equation}
Now let $\alpha\uparrow1$ in the interior point-process limit.  Equations
\eqref{eq:tail-tightness} and~\eqref{eq:intensity-tail-bound} give the full
convergence~\eqref{eq:point-process-limit}, and hence
$X_n\Rightarrow\Poi(\lambda)$.

For integer-valued random variables, convergence in distribution to a proper
integer-valued law implies pointwise convergence of the probability mass
functions; Scheffe's lemma then upgrades this to total variation convergence,
giving~\eqref{eq:count-poisson}.
\end{proof}

\section{Short cycles}\label{sec:short-cycles}

This section proves Theorem~\ref{thm:short-cycles}.  There are again two
tasks.  In the bulk, a linear reservoir of moderate-rate clocks yields the
finite-dimensional local law needed for factorial moments.  Near the
endpoint, the fixed-point estimate no longer suffices: the label that closes
a late cycle may come from anywhere.  We handle this by inserting every
proposed finite path into a common background permutation.  The construction
retains the essential permutation constraint that each rank has one
preimage.

For an integer $m$, write $[m]^r_{\ne}$ for the set of ordered $r$-tuples of
distinct elements of $[m]$.

\begin{lemma}\label{lem:moderate-reservoir}
Under Assumption~\ref{ass:profile}, for every $\alpha<1$ there are
$d,\eta>0$ such that, for all sufficiently large $n$,
\begin{equation}\label{eq:moderate-reservoir}
 \#\{i:\theta_{n,i}\ge d\}\ge(\alpha+2\eta)n.
\end{equation}
Consequently, after deleting a fixed number of clocks and before $\alpha n$
arrivals, the remaining total rate is at least $d\eta n$.
\end{lemma}

\begin{proof}
Since $f>0$ almost everywhere, choose $d,\eta>0$ so that
$|\{x:f(x)\ge2d\}|>\alpha+3\eta$.  The part of this set on which $f_n<d$ has
measure at most $d^{-1}\lVert f_n-f\rVert_1$, which proves
\eqref{eq:moderate-reservoir}.  At most $\alpha n+O(1)$ members of the
reservoir have arrived or been deleted, leaving at least $\eta n$ clocks of
rate at least $d$.
\end{proof}

\begin{lemma}\label{lem:cyclic-local}
Under Assumption~\ref{ass:profile}, fix $r\ge1$, $\alpha<1$, a permutation
$\tau\in S_r$, and a continuous function $g:[0,\alpha]^r\to\R$.  Then
\begin{align}
 &\sum_{\boldsymbol i\in[\lfloor\alpha n\rfloor]^r_{\ne}}
 g(\boldsymbol i/n)
 \Pp\bigl(R_{n,i_a}=i_{\tau(a)},\ 1\le a\le r\bigr)
 \label{eq:cyclic-local-limit}\\
 &\hspace{0.7in}\longrightarrow
 \int_{[0,\alpha]^r}g(\boldsymbol x)
 \prod_{a=1}^r\rho(x_a,x_{\tau(a)})\,d\boldsymbol x.
 \notag
\end{align}
Moreover, uniformly over distinct marked labels $i_1,\ldots,i_r$ and
distinct ranks $j_1,\ldots,j_r\le\alpha n$,
\begin{equation}\label{eq:weighted-bulk-cylinder}
 \Pp(R_{n,i_a}=j_a,\ 1\le a\le r)
 \le K_{r,\alpha}n^{-r}\prod_{a=1}^r\theta_{n,i_a}.
\end{equation}
\end{lemma}

\begin{proof}
Delete the marked clocks and denote the unmarked order statistics by
$T_0=0<T_1<\cdots<T_{n-r}$.  If $W_q^\circ$ is the remaining unmarked rate,
the memoryless property gives
\begin{equation}\label{eq:race-gap-representation}
 T_{q+1}-T_q=\frac{\xi_q}{W_q^\circ},
\end{equation}
where, conditionally on the elimination chain, the $\xi_q$ are independent
standard exponentials.  Lemma~\ref{lem:moderate-reservoir} gives
$W_q^\circ\ge d\eta n$ for $q\le\alpha n$.  Sorting the prescribed ranks,
and grouping marked clocks which must enter the same unmarked gap, therefore
gives
\[
 \Pp(R_{n,i_a}=j_a,\ a\le r)
 \le K_{r,\alpha}n^{-r}\prod_{a=1}^r\theta_{n,i_a},
\]
which is \eqref{eq:weighted-bulk-cylinder}.

We identify the asymptotic first with all marked rates bounded by a fixed
$M$.  Deleting finitely many such clocks changes the empirical race by
$o(1)$, uniformly in their labels.  Lemma~\ref{lem:race-law} and
\eqref{eq:race-gap-representation} then show, uniformly when the prescribed
ranks are separated by at least $\zeta n$, that
\begin{equation}\label{eq:bounded-marked-asymptotic}
 n^r\Pp(R_{n,i_a}=j_a,\ a\le r)
 -\prod_{a=1}^r
 \frac{\theta_{n,i_a}e^{-\theta_{n,i_a}t_{j_a/n}}}
 {D(t_{j_a/n})}
 \longrightarrow0.
\end{equation}
Indeed, each marked clock is inserted into a relevant unmarked gap; Taylor
expansion of its exponential density is uniform for rates at most $M$, and
the rescaled separated gaps converge jointly to independent standard
exponentials.

The contribution to \eqref{eq:cyclic-local-limit} of tuples containing a
rate larger than $M$ is, by \eqref{eq:weighted-bulk-cylinder}, at most a
constant times
\begin{equation}\label{eq:bulk-high-truncation}
 \sum_{a=1}^r
 \left(\frac1n\sum_i\theta_{n,i}\one\{\theta_{n,i}>M\}\right)
 \left(\frac1n\sum_i\theta_{n,i}\right)^{r-1}.
\end{equation}
It vanishes as first $n\to\infty$ and then $M\to\infty$.  After this rate
truncation, tuples with two coordinates within $\zeta n$ have contribution
$O_{r,M}(\zeta)$; the discarded high-rate part is still controlled by
\eqref{eq:bulk-high-truncation}.

It remains to pass from the grid kernel in
\eqref{eq:bounded-marked-asymptotic} to $\rho$.  On rates at most $M$, the
kernel is bounded and uniformly continuous in its rank coordinate on
$[0,\alpha]$, so the grid sum differs from the corresponding step integral by
$o_M(1)$.  The map $a\mapsto ae^{-at}$ is one-Lipschitz, uniformly for
$a,t\ge0$, and $D(t_y)$ is bounded below for $y\le\alpha$.  A telescoping
product and $\lVert f_n-f\rVert_1\to0$ identify the integral.  The parts on
which either $f_n$ or $f$ exceeds $M$ are bounded by their $L^1$ tails,
because every numerator is at most its source rate.  Let $n\to\infty$, then
$\zeta\downarrow0$, and finally $M\to\infty$ to obtain
\eqref{eq:cyclic-local-limit}.
\end{proof}

We now set up the endpoint resampling argument.  Fix a proposed $\ell$-cycle
and abbreviate $\theta_{n,i}$ to $\theta_i$.  Generate a background family
$E_i^0\sim\operatorname{Exp}(\theta_i)$, with order statistics
\[
 0=T_0<T_1<\cdots<T_n<T_{n+1}=\infty,
\]
and replace the clocks on the proposed cycle by independent copies.  Put
\begin{equation}\label{eq:ghost-window}
 J_j:=\bigl(T_{(j-\ell)\vee0},T_{(j+\ell)\wedge(n+1)}\bigr),
 \qquad
 p_{i,j}:=\int_{J_j}\theta_i e^{-\theta_i t}\,dt.
\end{equation}
The window $J_j$ contains every time that could acquire rank $j$ after these
$\ell$ replacements.  Since replacing $\ell$ clocks changes every rank by at
most $\ell$, averaging over the background gives
\begin{equation}\label{eq:ghost-cylinder-bound}
 \Pp(R_{i_a}=j_a,\ 1\le a\le\ell)
 \le\E\prod_{a=1}^\ell p_{i_a,j_a}.
\end{equation}
Each time belongs to at most $2\ell+1$ windows, and consequently
\begin{equation}\label{eq:ghost-row-bound}
 \sum_{j=1}^np_{i,j}\le2\ell+1.
\end{equation}

The next lemma iterates this insertion bound along a finite path.  Its
distinct-source condition is exactly the condition supplied by a genuine
cycle.

\begin{lemma}\label{lem:finite-insertion-path}
For fixed $\ell,m\ge1$ and every $v\in[n]$,
\begin{equation}\label{eq:finite-insertion-path}
 \E\sum_{\substack{u_0,\ldots,u_{m-1}\in[n]\\
                    u_0,\ldots,u_{m-1}\ {\mathrm{distinct}}}}
 p_{u_0,u_1}\cdots p_{u_{m-2},u_{m-1}}p_{u_{m-1},v}
 \le K_{\ell,m}.
\end{equation}
\end{lemma}

\begin{proof}
Generate one independent replacement $E_i^1$ for every background clock.
For a fixed source sequence $\boldsymbol u$, the product in
\eqref{eq:finite-insertion-path} is the conditional probability that every
$E_{u_a}^1$ falls in the window indexed by the next vertex, ending with
$J_v$.  Swap $E_{u_a}^0$ with $E_{u_a}^1$ for all $a$.  These swaps concern
distinct labels and preserve the joint law.  If $\widetilde R$ is the rank
permutation of the new background, success before the swap implies
\begin{equation}\label{eq:approximate-permutation-path}
 |\widetilde R_{u_a}-u_{a+1}|\le q_{\ell,m}\quad(a<m-1),
 \qquad
 |\widetilde R_{u_{m-1}}-v|\le q_{\ell,m},
\end{equation}
where $q_{\ell,m}=\ell+m+2$ is sufficient: membership in an old window costs
at most $\ell+1$ ranks, and the $m$ swaps cost at most another $m$.

For a deterministic permutation $\sigma$, give $u$ an edge to every $w$ with
$|\sigma(u)-w|\le q_{\ell,m}$.  Each target has indegree at most
$2q_{\ell,m}+1$, since every rank has one preimage.  Thus at most
$(2q_{\ell,m}+1)^m$ length-$m$ paths end at $v$.  Sum the
measure-preserving comparison over $\boldsymbol u$.
\end{proof}

Put $c_j=\sum_kp_{k,j}$.  Expanding this sum and separating whether its source
is new gives the useful consequence
\begin{equation}\label{eq:added-predecessor}
 \sup_{v\in[n]}\E
 \sum_{\substack{i_2,\ldots,i_\ell\in[n]\setminus\{v\}\\
                  i_2,\ldots,i_\ell\ {\mathrm{distinct}}}}
 c_{i_2}p_{i_2,i_3}\cdots p_{i_\ell,v}\le K_\ell.
\end{equation}
Indeed, when the source $k$ in $c_{i_2}=\sum_kp_{k,i_2}$ is distinct from
$i_2,\ldots,i_\ell$, Lemma~\ref{lem:finite-insertion-path} applies.  If
$k=i_a$, discard the extra factor $p_{i_a,i_2}\le1$ and apply the lemma to
the remaining path.  The interpretation for $\ell=1$ is
$\sup_v\E c_v<\infty$.

We first prepare two truncations: high rates globally and low rates on compact
label intervals.  For $M>1$ and $\delta>0$, let
\[
 H_{n,M}:=\{i:\theta_{n,i}>M\},\qquad
 L_{n,\delta}:=\{i:\theta_{n,i}<\delta\},
\]
and let $V_{n,L}(S)$ be the number of vertices in $S$ which belong to cycles
of length at most $L$.

\begin{proposition}[Rate truncation]\label{prop:cycle-rate-truncation}
Under Assumption~\ref{ass:profile},
\begin{align}
 \lim_{M\to\infty}\limsup_n\E V_{n,L}(H_{n,M})&=0,
 \label{eq:high-cycle-truncation}\\
 \lim_{\delta\downarrow0}\limsup_n
 \E V_{n,L}\bigl(L_{n,\delta}\cap[\lfloor\beta n\rfloor]\bigr)&=0
 \qquad(\beta<1).
 \label{eq:compact-low-cycle-truncation}
\end{align}
\end{proposition}

\begin{proof}
The $L^1$ profile convergence gives
\begin{equation}\label{eq:profile-tail-controls}
 \lim_{M\to\infty}\limsup_n\frac1n
 \sum_{i\in H_{n,M}}\theta_{n,i}=0,
 \qquad
 \lim_{\delta\downarrow0}\limsup_n\frac{|L_{n,\delta}|}{n}=0.
\end{equation}
For the second relation, compare $\{f_n<\delta\}$ with
$\{f<2\delta\}$ and use $f>0$ almost everywhere.

Choose $M$ so large that the labels of rate at most $M$ carry total weight at
least $n/2$.  If $\theta_i>M$, then, with
$G_n(t)=\sum_k\theta_k e^{-\theta_kt}$,
\[
 \theta_i e^{-\theta_i t}\le\frac{2\theta_i}{n}G_n(t).
\]
Indeed, all rates at most $M$ have survival factor at least
$e^{-\theta_i t}$.  Since $c_j=\int_{J_j}G_n(t)dt$, this gives
\begin{equation}\label{eq:high-entry-comparison}
 p_{i,j}\le\frac{2\theta_i}{n}c_j.
\end{equation}
Start a proposed cycle at $i\in H_{n,M}$, apply
\eqref{eq:ghost-cylinder-bound} and then
\eqref{eq:high-entry-comparison} to its first edge.  For $\ell\ge2$,
\begin{align*}
 \Pp(i\text{ belongs to an }\ell\text{-cycle})
 &\le \E\sum_{\substack{i_2,\ldots,i_\ell\ne i\\
                         i_2,\ldots,i_\ell\ {\mathrm{distinct}}}}
 p_{i,i_2}p_{i_2,i_3}\cdots p_{i_\ell,i}\\
 &\le \frac{2\theta_i}{n}\E
 \sum_{\substack{i_2,\ldots,i_\ell\ne i\\
                  i_2,\ldots,i_\ell\ {\mathrm{distinct}}}}
 c_{i_2}p_{i_2,i_3}\cdots p_{i_\ell,i}
 \le K_\ell\frac{\theta_i}{n}
\end{align*}
by \eqref{eq:added-predecessor}.  For fixed points, use
$p_{i,i}\le2\theta_i c_i/n$ and the $\ell=1$ interpretation of that estimate.
Summing over $i$ and $\ell\le L$ gives
\begin{equation}\label{eq:quantitative-high-cycles}
 \E V_{n,L}(H_{n,M})
 \le\frac{K_L}{n}\sum_{i\in H_{n,M}}\theta_{n,i}.
\end{equation}
This proves \eqref{eq:high-cycle-truncation}.  Now fix $\beta<1$.  Lemma
\ref{lem:moderate-reservoir}, applied at a fixed number larger than $\beta$,
and \eqref{eq:race-gap-representation} give
\begin{equation}\label{eq:interior-window-length}
 \sup_{i\le\beta n}\E|J_i|\le K_L/n.
\end{equation}
In a cycle avoiding $H_{n,M}$, the edge closing into $i$ is at most
$M|J_i|$.  Sum all other edges with \eqref{eq:ghost-row-bound} to get
\[
 \Pp\bigl(i\text{ lies in a short cycle avoiding }H_{n,M}\bigr)
 \le K_{L,M}/n.
\]
A cycle meeting both sets contains at most $L$ vertices from
$L_{n,\delta}\cap[\lfloor\beta n\rfloor]$ per high-rate vertex.  Therefore
\[
 \E V_{n,L}\bigl(L_{n,\delta}\cap[\lfloor\beta n\rfloor]\bigr)
 \le L\E V_{n,L}(H_{n,M})
   +K_{L,M}\frac{|L_{n,\delta}|}{n}.
\]
Use \eqref{eq:profile-tail-controls}, then let $M\to\infty$, to prove
\eqref{eq:compact-low-cycle-truncation}.
\end{proof}

For $\alpha<1$, let $C_{n,\ell}^{(\alpha)}$ count cycles all of whose labels
are at most $\alpha n$.

\begin{proposition}[Short-cycle shell tightness]
\label{prop:cycle-endpoint-bound}
Under Assumptions~\ref{ass:profile} and~\ref{ass:fixed-endpoint}, for every
fixed $L$,
\begin{equation}\label{eq:cycle-tail-tightness}
 \lim_{\alpha\uparrow1}\limsup_n
 \E\sum_{\ell=1}^L
 \bigl(C_{n,\ell}-C_{n,\ell}^{(\alpha)}\bigr)=0.
\end{equation}
\end{proposition}

\begin{proof}
The case $\ell=1$ is Proposition~\ref{prop:endpoint-bound}.  Fix
$2\le\ell\le L$ and put $B=2\ell+1$.  We first restrict to cycles avoiding
$H_{n,M}$.  Proposition~\ref{prop:cycle-rate-truncation} shows that the
discarded contribution vanishes as $M\to\infty$.

Fix a large shell cutoff $J_0$ and then $\beta<1$ so close to one that, for
all sufficiently large $n$, every label larger than $\beta n$ belongs to a
shell with index at least $J_0$.  We also discard cycles meeting
$L_{n,\delta}\cap[\lfloor\beta n\rfloor]$.  Their contribution vanishes as
$\delta\downarrow0$ by Proposition~\ref{prop:cycle-rate-truncation}.
Thus every remaining cycle vertex has rate at most $M$, and every remaining
vertex at most $\beta n$ has rate at least $\delta$.

Rotate a proposed cycle uniquely as
\[
 u\longrightarrow v\longrightarrow i_3\longrightarrow\cdots
 \longrightarrow i_\ell\longrightarrow u,
 \qquad v=\max\{u,v,i_3,\ldots,i_\ell\}.
\]
Conditional on the common background in \eqref{eq:ghost-window}, let
\[
 A_{v,u}:=\sum_{i_3,\ldots,i_\ell}
 p_{v,i_3}\cdots p_{i_\ell,u},
\]
where the intermediate sums range over retained labels; both distinctness
and the genuine-cycle restrictions $i_a<v$ are dropped, while the outer pair
remains restricted by $u<v$.  For $\ell=2$ this means
$A_{v,u}=p_{v,u}$.  Iterating
\eqref{eq:ghost-row-bound} gives the pointwise bound
\begin{equation}\label{eq:marked-return-bound}
 \sum_u A_{v,u}\le B^{\ell-1}.
\end{equation}
Consequently, for any set $S$ of retained predecessors and every $t$,
\begin{equation}\label{eq:marked-edge-bound}
 \sum_v\one\{t\in J_v\}\sum_{u\in S}
 \theta_u e^{-\theta_ut}A_{v,u}
 \le B^\ell\sup_{u\in S}\theta_u e^{-\theta_ut},
\end{equation}
because at most $B$ target windows contain $t$.
Here the supremum is zero when $S$ is empty, and enlarging the actual sum
from $u<v$ to all pairs can only increase its left side.
The unique rotation cancels the factor $1/\ell$ in the definition of the
cycle count.  Hence \eqref{eq:ghost-cylinder-bound} gives
\[
 \E\bigl[\#\{\text{retained $\ell$-cycles}\}\bigr]
 \le \E\sum_{u<v}p_{u,v}A_{v,u},
\]
with any further restrictions imposed on the outer pair in both sums.

Suppose first that the terminal depth of $v$ belongs to
$\mathcal D_{n,j}$, $j\ge J$.  The part of $p_{u,v}$ coming from
$t<j-h_j$ can be nonzero only if
\[
 T_{n-r_{n,j}+1-\ell}<j-h_j.
\]
Indeed, $v\ge n-r_{n,j}+1$ and $J_v$ starts at $T_{v-\ell}$.  Jensen's
inequality and the normalization give
\[
 S(j-h_j)\ge ne^{-j+h_j}\ge r_{n,j}e^{h_j}.
\]
For all large $j$, a Chernoff lower-tail bound therefore yields
\[
 \Pp\bigl(T_{n-r_{n,j}+1-\ell}<j-h_j\bigr)
 \le \exp\{-c r_{n,j}e^{h_j}\}.
\]
Pointwise in the common background, the early part of $p_{u,v}$ is at most
the indicator of this event, so the total rooted weight for each $v$ is at
most $B^{\ell-1}$ times that indicator.  Taking expectations and using
$|\mathcal D_{n,j}|\le r_{n,j}$, all early-window contributions are bounded
by a constant times
\begin{equation}\label{eq:early-cycle-shells}
 \sum_{\substack{j\ge J\\\mathcal D_{n,j}\ne\varnothing}}
 r_{n,j}e^{-c r_{n,j}e^{h_j}}
 \le\sum_{r\ge1}r e^{-c r e^{h_J}}
 \longrightarrow0.
\end{equation}
Here the maxima $r_{n,j}$ of the nonempty shells are distinct integers and
$h_j\ge h_J$ for $j\ge J$.

It remains to integrate $p_{u,v}$ over $t\ge j-h_j$.  If $u>\beta n$, its
terminal depth belongs to a shell $\mathcal D_{n,r}$ with $r\ge J_0$.
Since $u<v$, one has $r\le j$, and Lemma
\ref{lem:diagonal-shell-buffer} gives
$t\ge j-h_j\ge r-h_r$.  By increasing $J_0$, we may make the limsup in
\eqref{eq:buffered-shell-condition} smaller than $\mathcal Q(1)$; then, for
all sufficiently large $n$, every nonempty source shell $r\ge J_0$ satisfies
$b_{n,r}(r-h_r)>1$.  As
$b_{n,r}\le\theta_u\le M$, the map $a\mapsto ae^{-at}$ is decreasing on
this range, and hence
\[
 \theta_u e^{-\theta_u t}\le b_{n,r}e^{-b_{n,r}t}.
\]
Apply \eqref{eq:marked-edge-bound} separately to each source shell,
retaining $u<v$ (and hence $r\le j$) before enlarging the time range to
$[r-h_r,\infty)$, and then sum over the shells.  Their contribution is
at most
\begin{equation}\label{eq:late-cycle-shells}
 B^\ell\sum_{\substack{r\ge J_0\\\mathcal D_{n,r}\ne\varnothing}}
 \int_{r-h_r}^\infty b_{n,r}e^{-b_{n,r}t}\,dt
 =B^\ell\sum_{\substack{r\ge J_0\\\mathcal D_{n,r}\ne\varnothing}}
 e^{-b_{n,r}(r-h_r)},
\end{equation}
which vanishes as $J_0\to\infty$ by
\eqref{eq:buffered-shell-condition}.

Finally, if $u\le\beta n$, then $\delta\le\theta_u\le M$.  For $J$ large,
$\delta(J-h_J)>1$, so on $t\ge j-h_j\ge J-h_J$,
\[
 \theta_u e^{-\theta_u t}\le\delta e^{-\delta t}.
\]
The same marked-edge bound makes this contribution at most
\begin{equation}\label{eq:late-cycle-bulk}
 B^\ell\int_{J-h_J}^\infty\delta e^{-\delta t}\,dt
 =B^\ell e^{-\delta(J-h_J)},
\end{equation}
which vanishes as $J\to\infty$.

Choose first $M$, then $J_0$ and the corresponding $\beta$, then $\delta$,
and finally $J$, keeping these cutoffs fixed before taking each limsup in
$n$.  Equations
\eqref{eq:early-cycle-shells}--\eqref{eq:late-cycle-bulk}, followed by the
two truncation estimates, prove \eqref{eq:cycle-tail-tightness}; replacing
$1-e^{-J}$ by an arbitrary $\alpha\uparrow1$ only requires a rounding check.
Integer rounding can add at most one terminal-depth label, which for all
large $n$ lies in a shell with index at least $J-1$.  Thus the estimate with
$J-1$ proves the result along $\alpha=1-e^{-J}$, and monotonicity gives the
general limit.
\end{proof}

\begin{proof}[Proof of Theorem~\ref{thm:short-cycles}]
Fix $\alpha<1$ and nonnegative integers $m_1,\ldots,m_L$, and put
$r=\sum_{\ell=1}^L\ell m_\ell$.  Expanding the joint falling factorial gives
an ordered collection of vertex-disjoint cycles.  Lemma
\ref{lem:cyclic-local} therefore yields
\begin{equation}\label{eq:cycle-factorial-moments}
 \E\prod_{\ell=1}^L(C_{n,\ell}^{(\alpha)})_{m_\ell}
 \longrightarrow
 \prod_{\ell=1}^L\lambda_\ell(\alpha)^{m_\ell},
\end{equation}
where
\[
 \lambda_\ell(\alpha):=\frac1\ell
 \int_{[0,\alpha]^\ell}
 \prod_{a=1}^\ell\rho(x_a,x_{a+1})\,d\boldsymbol x,
 \qquad x_{\ell+1}=x_1.
\]
These are the joint factorial moments of independent Poisson variables, so
the truncated cycle vector has the asserted limit.

Proposition~\ref{prop:cycle-endpoint-bound} gives the full endpoint
tightness~\eqref{eq:cycle-tail-tightness}.

First-moment convergence and \eqref{eq:cycle-tail-tightness} show that
\[
 \sup_{\alpha<\beta<1}
 \bigl[\lambda_\ell(\beta)-\lambda_\ell(\alpha)\bigr]
 \longrightarrow0
 \qquad(\alpha\uparrow1).
\]
Hence the increasing $\lambda_\ell(\alpha)$ have finite
limits, necessarily the integrals in \eqref{eq:cycle-intensity}.  Markov's
inequality removes the label cutoff in \eqref{eq:cycle-factorial-moments} and
proves \eqref{eq:cycle-vector-limit}.  Finally, weak convergence on the
countable lattice $\mathbb Z_+^L$ gives pointwise convergence of mass
functions; Scheffe's lemma upgrades it to total variation convergence.
\end{proof}

\section{Power-law endpoint singularities and cycle CLTs}
\label{sec:power-law}

\subsection{Triangular-array hypotheses}

We now prove the endpoint theorems in a form stronger than the profile
statement in the introduction.  Two inputs must be distinguished.  The local
power law determines the rate of a label at a given distance from a corner;
the global population law determines how many clocks have comparable rates.
The corner is active only when these two scales match.

The midpoint and endpoint sampling conventions differ by a bounded shift of
the endpoint grid.  To cover both, fix $\delta_0,\delta_1\in[0,1)$ and put
\[
 s_{n,m}^{(j)}:=(m-\delta_j)/n,
 \qquad j\in\{0,1\}.
\]

\begin{assumption}\label{ass:power-endpoints}
There are $\eps_0,\eta,c_0,c_1>0$ and $r_0,r_1\in\R$ such that, uniformly for
$m\le \eps_0n$,
\begin{align}
 \theta_{n,m}&=c_0(s_{n,m}^{(0)})^{r_0}
   (1+O((s_{n,m}^{(0)})^{\eta})),\label{eq:left-endpoint-power}\\
 \theta_{n,n-m+1}&=c_1(s_{n,m}^{(1)})^{r_1}
   (1+O((s_{n,m}^{(1)})^{\eta})).\label{eq:right-endpoint-power}
\end{align}
\end{assumption}

Define the empirical transforms
\begin{equation}\label{eq:population-transforms}
 H_n(t):=\frac1n\sum_{i=1}^ne^{-t\theta_{n,i}},\qquad
 G_n(t):=\frac1n\sum_{i=1}^n(1-e^{-t\theta_{n,i}}),\qquad
 D_n(t):=\frac1n\sum_{i=1}^n\theta_{n,i}e^{-t\theta_{n,i}}.
\end{equation}
Thus $H_n$ measures clocks still running at a large time, whereas $G_n$
measures clocks that have rung by a small time.

\begin{assumption}\label{ass:power-population}
Whenever $r_1>0$, there are $\beta_{\rm s}\ge r_1$, $A_->0$, and $\xi_->0$
such that, uniformly as $t\to\infty$ with
$nt^{-1/\beta_{\rm s}}\to\infty$,
\begin{align}
 H_n(t)&=A_-t^{-1/\beta_{\rm s}}
 \left(1+O\left(t^{-\xi_-}
 +(nt^{-1/\beta_{\rm s}})^{-1}\right)\right),
 \label{eq:power-discrete-laplace}\\
 D_n(t)&=\frac{A_-}{\beta_{\rm s}}t^{-1/\beta_{\rm s}-1}
 \left(1+O\left(t^{-\xi_-}
 +(nt^{-1/\beta_{\rm s}})^{-1}\right)\right).
 \label{eq:power-discrete-weighted-laplace}
\end{align}
Whenever $r_0<-1$, there are $\alpha_{\rm f}\ge-r_0$, $A_+>0$, and
$\xi_+>0$ such that, uniformly as $t\downarrow0$ with
$nt^{1/\alpha_{\rm f}}\to\infty$,
\begin{align}
 G_n(t)&=A_+t^{1/\alpha_{\rm f}}
 \left(1+O\left(t^{\xi_+}
 +(nt^{1/\alpha_{\rm f}})^{-1}\right)\right),
 \label{eq:power-discrete-arrivals}\\
 D_n(t)&=\frac{A_+}{\alpha_{\rm f}}t^{1/\alpha_{\rm f}-1}
 \left(1+O\left(t^{\xi_+}
 +(nt^{1/\alpha_{\rm f}})^{-1}\right)\right).
 \label{eq:power-discrete-arrival-weight}
\end{align}
At least one matching relation $\beta_{\rm s}=r_1>0$ or
$\alpha_{\rm f}=-r_0>1$ holds.
\end{assumption}

The right corner is active when $\beta_{\rm s}=r_1=:\beta>0$, and the left
corner is active when $\alpha_{\rm f}=-r_0=:\alpha>1$.  At active corners
define
\begin{align}
 q_R&:=c_1A_-^\beta,
 &b_{R,\ell}&:=\mathfrak b_{\beta,\ell}(q_R),
 \label{eq:general-q-right}\\
 q_L&:=c_0A_+^{-\alpha},
 &b_{L,\ell}&:=\mathfrak b_{\alpha,\ell}(q_L).
 \label{eq:general-q-left}
\end{align}
Set the coefficient at an inactive corner to zero and again write
$B_\ell=b_{R,\ell}+b_{L,\ell}$.  Under
Assumptions~\ref{ass:power-endpoints} and~\ref{ass:power-population},
the conclusions of Theorems~\ref{thm:power-clt}
and~\ref{thm:power-spatial} hold with these constants.

The regular-variation assumption has a direct interpretation for a sampled
profile.  If
\begin{align*}
 |\{x:f(x)\le y\}|&\sim K_-y^{1/\beta_{\rm s}}
 &&(y\downarrow0),\\
 |\{x:f(x)>y\}|&\sim K_+y^{-1/\alpha_{\rm f}}
 &&(y\to\infty),
\end{align*}
then
\[
 A_-=K_-\Gamma(1+1/\beta_{\rm s}),\qquad
 A_+=K_+\Gamma(1-1/\alpha_{\rm f}).
\]
Thus same-order singularities anywhere in the profile add to $A_-$ or
$A_+$; lower-order ones disappear from the leading constant.  A stronger
unaligned singularity makes the corresponding inequality strict and turns
that corner inactive.

A convenient sufficient condition makes precise when the endpoints alone
determine the global exponent.  Put
\[
 I_n:=\{i:\eps_0n<i<(1-\eps_0)n\},\quad
 H_n^I(t):=\frac1n\sum_{i\in I_n}e^{-t\theta_{n,i}},\quad
 G_n^I(t):=\frac1n\sum_{i\in I_n}(1-e^{-t\theta_{n,i}}).
\]
For a right zero of order $\beta=r_1$, suppose no endpoint zero has order
larger than $\beta$ and, for some $\delta_->0$,
\begin{equation}\label{eq:interior-slow-negligible}
 H_n^I(t)=O(t^{-1/\beta-\delta_-}+n^{-1}),\qquad t\ge1.
\end{equation}
Then the slow-clock part of Assumption~\ref{ass:power-population} holds with
$\beta_{\rm s}=\beta$ and
\begin{equation}\label{eq:general-A-minus}
 A_-^{\rm end}=\Gamma(1+1/\beta)
 \sum_{j:\,r_j=\beta}c_j^{-1/\beta}.
\end{equation}
Likewise, for a left pole of order $\alpha=-r_0>1$, if no endpoint pole has
larger order and, for some $\delta_+>0$,
\begin{equation}\label{eq:interior-fast-negligible}
 G_n^I(t)=O(t^{1/\alpha+\delta_+}+n^{-1}),\qquad 0<t\le1,
\end{equation}
then the fast-clock part holds with $\alpha_{\rm f}=\alpha$ and
\begin{equation}\label{eq:general-A-plus}
 A_+^{\rm end}=\Gamma(1-1/\alpha)
 \sum_{j:\,-r_j=\alpha}c_j^{1/\alpha}.
\end{equation}
The two bounds follow, for example, from
\begin{equation}\label{eq:interior-moment-sufficient}
 \sup_n\frac1n\sum_{i\in I_n}\theta_{n,i}^{-p_-}<\infty
 \quad\text{for some }p_->1/\beta,
 \qquad
 \sup_n\frac1n\sum_{i\in I_n}\theta_{n,i}^{p_+}<\infty
 \quad\text{for some }1/\alpha<p_+\le1,
\end{equation}
when the corresponding endpoint is under consideration.  These moment
bounds are only convenient sufficient conditions.  If the interior itself
has a matching regularly varying population, its constant is simply added
to $A_-^{\rm end}$ or $A_+^{\rm end}$.

The profile in Assumption~\ref{ass:simple-power-profile} satisfies these
stronger hypotheses with $\delta_0=\delta_1=1/2$.  At a right power zero,
$A_-=\Gamma(1+1/\beta)c_R^{-1/\beta}$; at a left power pole,
$A_+=\Gamma(1-1/\alpha)c_L^{1/\alpha}$.  Hence
\eqref{eq:general-q-right}--\eqref{eq:general-q-left} reduce exactly to
\eqref{eq:simple-q-right}--\eqref{eq:simple-q-left}.

For the critical exponent $r_0=-1$, the general version assumes that, for
some $C_+>0$, uniformly as $t\downarrow0$ with $nt\to\infty$,
\begin{align}
 G_n(t)&=C_+t\log(1/t)
 \left(1+O\left(\frac1{\log(1/t)}+(nt)^{-1}\right)\right),
 \label{eq:critical-arrival-tail}\\
 D_n(t)&=C_+\log(1/t)
 \left(1+O\left(\frac1{\log(1/t)}+(nt)^{-1}\right)\right).
 \label{eq:critical-arrival-weight}
\end{align}
Together with Assumption~\ref{ass:power-endpoints}, $r_0=-1$, and
$-1\le r_1\le0$, these conditions imply the conclusion of
Theorem~\ref{thm:critical-pole} in the more general form
\begin{equation}\label{eq:critical-pole-coefficient}
 \frac{C_{n,1}-d\log\log n}{\sqrt{d\log\log n}}
 \Longrightarrow\mathcal N(0,1),\qquad
 \E C_{n,1}\sim d\log\log n,\qquad d:=c_0/C_+.
\end{equation}
For the simple profile of Theorem~\ref{thm:critical-pole}, one has
$C_+=c_L$ and therefore $d=1$.

\subsection{One-point estimates at the two corners}

We first treat fixed points.  The upper-right corner is governed by slow
clocks that survive to large times; the lower-left corner is governed by fast
clocks that arrive at small times.  We estimate their predictable
compensators separately and add them at the end.  For $r\ge0$, put
\begin{equation}\label{eq:tail-moments}
 D_{n,r}(t):=\frac1n\sum_{i=1}^n
 \theta_{n,i}^r e^{-\theta_{n,i}t}.
\end{equation}
Thus $D_{n,1}=D_n$ from~\eqref{eq:population-transforms}.  We first prove the
sufficient conditions for the endpoint populations to determine $A_-$ and
$A_+$.  This proposition is not needed when a same-order intermediate
population has already been incorporated into either constant.

\begin{proposition}\label{prop:global-endpoint-tails}
Under Assumption~\ref{ass:power-endpoints}, the following hold.
\begin{enumerate}[label=(\roman*)]
\item If $\beta=r_1>0$, no endpoint zero has order larger than $\beta$,
and~\eqref{eq:interior-slow-negligible} holds, then the slow-clock part of
Assumption~\ref{ass:power-population} holds with
$\beta_{\rm s}=\beta$ and
$A_-=A_-^{\rm end}$ from~\eqref{eq:general-A-minus}.
\item If $\alpha=-r_0>1$, no endpoint pole has order larger than $\alpha$,
and~\eqref{eq:interior-fast-negligible} holds, then the fast-clock part of
Assumption~\ref{ass:power-population} holds with
$\alpha_{\rm f}=\alpha$ and
$A_+=A_+^{\rm end}$ from~\eqref{eq:general-A-plus}.
\end{enumerate}
\end{proposition}

\begin{proof}
Consider first a zero branch of exponent $\beta$ and coefficient $c$.
Write $s_m=(m-\delta)/n$ for its endpoint grid.
The functions $e^{-cts^\beta}$ and
$cs^\beta e^{-cts^\beta}$ have total variations $O(1)$ and $O(t^{-1})$.
Shifted Riemann summation and the change of variables
$s=t^{-1/\beta}z$ give
\begin{align*}
 \frac1n\sum_{m\le\eps_0n}e^{-ct s_m^\beta}
 &=\Gamma(1+1/\beta)c^{-1/\beta}t^{-1/\beta}
   +O(n^{-1}+e^{-c't}),\\
 \frac1n\sum_{m\le\eps_0n}cs_m^\beta e^{-ct s_m^\beta}
 &=\frac{\Gamma(1+1/\beta)c^{-1/\beta}}\beta
 t^{-1/\beta-1}+O((nt)^{-1}+e^{-c't}).
\end{align*}
The relative endpoint error in
\eqref{eq:left-endpoint-power}--\eqref{eq:right-endpoint-power} changes
these estimates by a power of $t^{-1}$.  The hypotheses in part~(i) say that
only branches of exponent $\beta$ contribute at leading order, and their
constants add to $A_-^{\rm end}$.  Equation
\eqref{eq:interior-slow-negligible} gives the required power saving for
$H_n^I(t)$, and
\[
 \frac1n\sum_{i\in I_n}\theta_{n,i}e^{-t\theta_{n,i}}
 \le \frac Ct H_n^I(t/2)
\]
gives the same relative saving for the weighted transform.  The additive
$n^{-1}$ error becomes $O((nt^{-1/\beta})^{-1})$ relative to the main
term.  This proves part~(i).

For a pole branch of exponent $\alpha>1$ and coefficient $c$, the functions
$1-e^{-ct s^{-\alpha}}$ and
$cs^{-\alpha}e^{-ct s^{-\alpha}}$ have total variations $O(1)$ and
$O(t^{-1})$.  Write $s_m=(m-\delta)/n$ for this endpoint grid.  Shifted
Riemann summation and $s=t^{1/\alpha}z$ give
\begin{align*}
 \frac1n\sum_{m\le\eps_0n}(1-e^{-ct s_m^{-\alpha}})
 &=\Gamma(1-1/\alpha)c^{1/\alpha}t^{1/\alpha}
   +O(n^{-1}+t),\\
 \frac1n\sum_{m\le\eps_0n}cs_m^{-\alpha}e^{-ct s_m^{-\alpha}}
 &=\frac{\Gamma(1-1/\alpha)c^{1/\alpha}}\alpha
 t^{1/\alpha-1}+O((nt)^{-1}+1).
\end{align*}
Only branches of exponent $\alpha$ contribute at leading order.
Equation~\eqref{eq:interior-fast-negligible} gives the corresponding power
saving for $G_n^I(t)$.  The inequality
\[
 \frac1n\sum_{i\in I_n}\theta_{n,i}e^{-t\theta_{n,i}}
 \le \frac Ct G_n^I(2t)
\]
does the same for the weighted transform.  The $n^{-1}$ term again matches
the discrete error in Assumption~\ref{ass:power-population}.  Endpoint errors
and subdominant pole branches save a positive power as before.  This proves
part~(ii).
\end{proof}

Let $k_m=n-m+1$.  At an active right corner write
$\beta:=r_1$, $A:=A_-$, $q:=q_R$, and define $t_{n,m}^-$ by
\begin{equation}\label{eq:tail-quantile}
 H_n(t_{n,m}^-)=m/n.
\end{equation}

\begin{proposition}
\label{prop:deterministic-slow}
If the right corner is active, let $L_n\to\infty$
and $\delta_n\downarrow0$, with $L_n\le\delta_n n$ eventually.  Uniformly for
$L_n\le m\le\delta_n n$,
\begin{align}
 t_{n,m}^-&=\left(\frac{An}{m}\right)^\beta
 \left(1+O((m/n)^{\xi}+m^{-1})\right),\label{eq:tm-power}\\
 \theta_{n,k_m}t_{n,m}^-&=q
 \left(1+O((m/n)^{\xi}+m^{-1})\right),\label{eq:theta-t-power}\\
 \frac{\theta_{n,k_m}}{nD_n(t_{n,m}^-)}
 &=\frac{\beta q}{m}
 \left(1+O((m/n)^{\xi}+m^{-1})\right)
 \label{eq:hazard-power}
\end{align}
for some $\xi>0$.
\end{proposition}

\begin{proof}
Invert \eqref{eq:power-discrete-laplace} at $H_n(t)=m/n$ and then use
\eqref{eq:right-endpoint-power}.  The error $t^{-\xi_-}$ becomes a positive
power of $m/n$, while replacing the shifted endpoint grid by $m/n$ costs
$O(m^{-1})$.  Finally use
\eqref{eq:power-discrete-weighted-laplace}.  Decreasing $\xi>0$ absorbs all
three error exponents.
\end{proof}

At an active left corner write $\alpha:=-r_0$, $A:=A_+$, $q:=q_L$,
and define $t_{n,m}^+$ by
\begin{equation}\label{eq:head-quantile}
 G_n(t_{n,m}^+)=m/n.
\end{equation}

\begin{proposition}
\label{prop:deterministic-early}
If the left corner is active, uniformly for
$L_n\le m\le\delta_n n$ as above,
\begin{align}
 t_{n,m}^+&=\left(\frac{m}{An}\right)^\alpha
 \left(1+O((m/n)^\xi+m^{-1})\right),
 \label{eq:early-time-power}\\
 \theta_{n,m}t_{n,m}^+&=q
 \left(1+O((m/n)^\xi+m^{-1})\right),
 \label{eq:early-theta-time}\\
 \frac{\theta_{n,m}}{nD_n(t_{n,m}^+)}
 &=\frac{\alpha q}{m}
 \left(1+O((m/n)^\xi+m^{-1})\right)
 \label{eq:early-hazard-power}
\end{align}
for some $\xi>0$.
\end{proposition}

\begin{proof}
Invert \eqref{eq:power-discrete-arrivals} at $G_n(t)=m/n$, substitute
\eqref{eq:left-endpoint-power}, and use
\eqref{eq:power-discrete-arrival-weight}.  The errors transform exactly as
in Proposition~\ref{prop:deterministic-slow}.
\end{proof}

The population estimates enter the fixed-point CLT only through the
predictable compensator.  We record the martingale step separately.  Recall
$p_{n,k}=\E(I_{n,k}\mid\cF_{n,k-1})$.

\begin{proposition}\label{prop:conditional-clt}
Suppose $v_n\to\infty$ and
\begin{equation}\label{eq:conditional-clt-input}
 \Lambda_n:=\sum_{k=1}^np_{n,k}
 =v_n+o_{\Pp}(\sqrt{v_n}),
 \qquad
 \sum_{k=1}^np_{n,k}^2=o_{\Pp}(v_n).
\end{equation}
Then
\[
 \frac{X_n-v_n}{\sqrt{v_n}}
 \ \Longrightarrow\ \mathcal N(0,1).
\]
If in addition $\E\Lambda_n\sim v_n$, then $\E X_n\sim v_n$.
\end{proposition}

\begin{proof}
Write
\begin{equation}\label{eq:martingale-decomp}
 X_n=M_n+\Lambda_n,\qquad
 M_n:=\sum_{k=1}^n(I_{n,k}-p_{n,k}).
\end{equation}
The predictable quadratic variation is
\begin{equation}\label{eq:quadratic-variation}
 V_n:=\sum_{k=1}^np_{n,k}(1-p_{n,k})
 =\Lambda_n-\sum_{k=1}^np_{n,k}^2.
\end{equation}
Thus $V_n/v_n\to1$ in probability.  The increments of $M_n$ are
bounded by one, so conditional Lindeberg is automatic after division by
$\sqrt{v_n}$.  The martingale central limit theorem
\cite{hall1980martingale} proves the claim, and $\E X_n=\E\Lambda_n$ proves
the expectation statement.
\end{proof}

Write $S_n(t):=nH_n(t)$.

\begin{lemma}
\label{lem:rough-slow}
There are $\delta_0,c_0,C_0>0$ such that, uniformly whenever
$m\to\infty$, $m\le\delta_0n$, $t=t_{n,m}^-$, and $t/2\le s\le2t$,
\begin{equation}\label{eq:rough-slow}
 t\asymp(n/m)^\beta,\qquad
 \theta_{n,k_m}t\asymp1,\qquad
 S_n(s)\asymp m,\qquad
 nD_{n,1}(s)\asymp\frac mt,\qquad
 nD_{n,2}(s)\le C_0\frac m{t^2}.
\end{equation}
Moreover, for $0<u\le1/2$,
\begin{align}
 S_n((1-u)t)&\ge m(1+c_0u),&
 S_n((1+u)t)&\le m(1-c_0u),
 \label{eq:count-mean-separation}\\
 \left|\frac{D_{n,1}((1\pm u)t)}{D_{n,1}(t)}-1\right|
 &\le C_0u.&&\label{eq:weight-mean-stability}
\end{align}
\end{lemma}

\begin{proof}
The discrete Laplace estimates
\eqref{eq:power-discrete-laplace}--\eqref{eq:power-discrete-weighted-laplace},
used with fixed multiplicative changes of $t$, give the first four comparisons
in \eqref{eq:rough-slow}.  Also,
\[
 D_{n,2}(s)\le\frac{C}{s}D_{n,1}(s/2),
\]
because $ae^{-as/2}\le C/s$; this proves the last comparison.  The logarithmic
derivatives satisfy
\[
 -\frac{d\log S_n(t)}{d\log t}
 =\frac{tnD_{n,1}(t)}{S_n(t)}\asymp1,
 \qquad
 -\frac{d\log D_{n,1}(t)}{d\log t}
 =\frac{tD_{n,2}(t)}{D_{n,1}(t)}\le C.
\]
Integrating these identities between $t$ and $(1\pm u)t$ proves
\eqref{eq:count-mean-separation} and \eqref{eq:weight-mean-stability}, after
decreasing $\delta_0$ and increasing the fixed lower threshold for $m$.
\end{proof}

Let $T_{n,m}:=\tau_{n,k_m}$, let
\begin{equation}\label{eq:random-populations}
 N_n(t):=\sum_i\one\{E_{n,i}\ge t\},\qquad
 W_n(t):=\sum_i\theta_{n,i}\one\{E_{n,i}\ge t\},
\end{equation}
and put
\begin{equation}\label{eq:slow-r-B}
 r_{n,m}:=\frac{\theta_{n,k_m}}{nD_n(t_{n,m}^-)},\qquad
 B_{n,m}:=\one\{E_{n,k_m}\ge t_{n,m}^-\}.
\end{equation}
The exact predictable probability is
\begin{equation}\label{eq:exact-tail-p}
 p_{n,k_m}
 =\frac{\theta_{n,k_m}\one\{E_{n,k_m}\ge T_{n,m}\}}
 {W_n(T_{n,m})}.
\end{equation}

\begin{proposition}
\label{prop:slow-compensator}
If $L_n\to\infty$, $\delta_n\downarrow0$, and
$L_n\le\delta_n n$, then
\begin{align}
 \E\sum_{m=L_n}^{\lfloor\delta_n n\rfloor}
 \left|p_{n,k_m}-r_{n,m}B_{n,m}\right|&\longrightarrow0,
 \label{eq:slow-L1}\\
 \E\sum_{m=L_n}^{\lfloor\delta_n n\rfloor}p_{n,k_m}^2
 &\longrightarrow0.\label{eq:slow-squares}
\end{align}
\end{proposition}

\begin{proof}
Fix $m$ in the stated range, abbreviate $t=t_{n,m}^-$, and set
$u=m^{-1/4}$ and $s_\pm=(1\pm u)t$.  Since $T_{n,m}>s_+$ forces
$N_n(s_+)\ge m$, while $T_{n,m}<s_-$ forces $N_n(s_-)\le m-1$,
the Poisson-binomial Chernoff bound and
\eqref{eq:count-mean-separation} give
\begin{equation}\label{eq:random-time-concentration}
 \Pp(T_{n,m}\notin[s_-,s_+])\le2e^{-c\sqrt m}.
\end{equation}

We next concentrate $W_n(s_-)$ and $W_n(s_+)$.  Truncate at
$b_m=K\log(m)/t$, with a fixed sufficiently large $K$.  For either sign, the
same fixed-factor estimates used in Lemma~\ref{lem:rough-slow} give
\begin{align*}
 \sum_{\theta_{n,i}>b_m}\Pp(E_{n,i}\ge s_\pm)
 &\le m^{-K/4}S_n(t/4)\le C m^{1-K/4},\\
 \sum_{\theta_{n,i}>b_m}\theta_{n,i}
       \Pp(E_{n,i}\ge s_\pm)
 &\le m^{-K/4}nD_n(t/4)\le C\frac mtm^{-K/4}.
\end{align*}
Thus, except on an event of probability $O(m^{1-K/4})$, no rate above $b_m$
survives to $s_\pm$; the omitted part of the mean is
$O(m^{1-K/4}/t)$.  The remaining summands are bounded by $b_m$ and their
total variance is at most $nD_{n,2}(s_\pm)\le Cm/t^2$.  Bernstein's
inequality, with deviation $u m/t$, therefore yields
\begin{equation}\label{eq:random-weight-concentration}
 \Pp\left(\left|W_n(s_\pm)-nD_n(s_\pm)\right|>C u\frac mt\right)
 \le C m^{-4}+Ce^{-c\sqrt m},
\end{equation}
after taking, say, $K=20$.  By monotonicity of $W_n$ and
\eqref{eq:weight-mean-stability}, equations
\eqref{eq:random-time-concentration}--\eqref{eq:random-weight-concentration}
imply that, outside an event of probability
\begin{equation}\label{eq:slow-bad-probability}
 b_m^\star:=Cm^{-4}+Ce^{-c\sqrt m},
\end{equation}
one has
\begin{equation}\label{eq:random-weight-at-order}
 T_{n,m}\in[s_-,s_+],\qquad
 W_n(T_{n,m})=nD_n(t)(1+O(u)).
\end{equation}

By Lemma~\ref{lem:rough-slow}, $r_{n,m}\le C/m$.  On the good event in
\eqref{eq:random-weight-at-order}, replacing the denominator in
\eqref{eq:exact-tail-p} costs at most $Cr_{n,m}u$.  The two survival
indicators can differ only if $E_{n,k_m}\in[s_-,s_+]$, an event of probability
at most $Cu$ because $\theta_{n,k_m}t\asymp1$.  On the bad event use
$p_{n,k_m}\le1$.  It follows that
\begin{align}
 \E|p_{n,k_m}-r_{n,m}B_{n,m}|
 &\le Cm^{-5/4}+Cb_m^\star,\label{eq:one-p-error}\\
 \E p_{n,k_m}^2
 &\le Cm^{-2}+Cb_m^\star.\label{eq:one-p-square}
\end{align}
Both right sides are summable from $m=L_n$ to infinity and their sums tend to
zero as $L_n\to\infty$.
\end{proof}

The fast corner has the same population estimate at small clock times.  Put
\begin{equation}\label{eq:early-r-B}
 \widehat r_{n,m}:=
 \frac{\theta_{n,m}}{nD_n(t_{n,m}^+)},
 \qquad
 \widehat B_{n,m}:=\one\{E_{n,m}\ge t_{n,m}^+\},
 \qquad
 \widehat T_{n,m}:=\tau_{n,m}.
\end{equation}

\begin{lemma}
\label{lem:rough-early}
If the left corner is active, there are $\delta_0,c,C>0$ such that,
uniformly when $m\to\infty$, $m\le\delta_0n$,
$t=t_{n,m}^+$, and $t/2\le s\le2t$,
\begin{equation}\label{eq:rough-early}
 t\asymp(m/n)^\alpha,\qquad
 \theta_{n,m}t\asymp1,\qquad
 nG_n(s)\asymp m,\qquad
 nD_{n,1}(s)\asymp\frac mt,\qquad
 nD_{n,2}(s)\le C\frac m{t^2}.
\end{equation}
Moreover, for $0<u\le1/2$,
\begin{align}
 nG_n((1-u)t)&\le m(1-cu),&
 nG_n((1+u)t)&\ge m(1+cu),
 \label{eq:early-count-separation}\\
 \left|\frac{D_{n,1}((1\pm u)t)}{D_{n,1}(t)}-1\right|
 &\le Cu.&&\label{eq:early-weight-stability}
\end{align}
\end{lemma}

\begin{proof}
Use \eqref{eq:power-discrete-arrivals}--
\eqref{eq:power-discrete-arrival-weight} at fixed multiplicative changes of
$t$.  The bound on $D_{n,2}$ and the logarithmic derivative argument are
identical to Lemma~\ref{lem:rough-slow}; for $G_n$, use
$tG_n'(t)/G_n(t)\asymp1$ rather than the logarithmic derivative of $H_n$.
\end{proof}

\begin{proposition}
\label{prop:early-compensator}
If the left corner is active and $L_n\to\infty$, $\delta_n\downarrow0$,
$L_n\le\delta_n n$, then
\begin{align}
 \E\sum_{m=L_n}^{\lfloor\delta_n n\rfloor}
 \left|p_{n,m}-\widehat r_{n,m}\widehat B_{n,m}\right|&\longrightarrow0,
 \label{eq:fast-L1}\\
 \E\sum_{m=L_n}^{\lfloor\delta_n n\rfloor}p_{n,m}^2
 &\longrightarrow0.\label{eq:fast-squares}
\end{align}
\end{proposition}

\begin{proof}
Fix $m$, put $t=t_{n,m}^+$, $u=m^{-1/4}$, and
$s_\pm=(1\pm u)t$.  The Poisson--binomial Chernoff bound and
\eqref{eq:early-count-separation} give
\[
 \Pp(\widehat T_{n,m}\notin[s_-,s_+])\le2e^{-c\sqrt m}.
\]
The concentration proof for $W_n(s_\pm)$ in
Proposition~\ref{prop:slow-compensator} uses only the fixed-factor
population estimates, the inequality
$D_{n,2}(s)\le Cs^{-1}D_{n,1}(s/2)$, and truncation at
$K\log(m)/t$.  Lemma~\ref{lem:rough-early} supplies the same inputs.
Consequently, outside an event of probability
$Cm^{-4}+Ce^{-c\sqrt m}$,
\[
 W_n(\widehat T_{n,m})=nD_n(t)(1+O(m^{-1/4})).
\]
Since $\widehat r_{n,m}\le C/m$ and
$\theta_{n,m}t\asymp1$, the last paragraph of the slow-corner proof gives
\begin{align*}
 \E|p_{n,m}-\widehat r_{n,m}\widehat B_{n,m}|
 &\le Cm^{-5/4}+Cm^{-4}+Ce^{-c\sqrt m},\\
 \E p_{n,m}^2&\le Cm^{-2}+Cm^{-4}+Ce^{-c\sqrt m}.
\end{align*}
Summation proves the proposition.
\end{proof}

\subsection{The critical pole}

At exponent one, the fast-clock population acquires a slowly varying
logarithmic factor.  The same compensator argument still applies, but the
mean grows like $\log\log n$ rather than $\log n$.

\begin{proof}[Proof of Theorem~\ref{thm:critical-pole}]
We use the global population laws
\eqref{eq:critical-arrival-tail}--\eqref{eq:critical-arrival-weight}.
Shifted Riemann summation shows that a displayed branch
$cs^{-1}$ contributes $ct\log(1/t)+O(t+n^{-1})$ to $G_n$ and
$c\log(1/t)+O(1+(nt)^{-1})$ to $D_n$.  Thus, in the endpoint-dominated
case, $C_+=\sum_{j:r_j=-1}c_j$; same-order intermediate contributions are
already included in the assumed value of $C_+$.

Let $t^0_{n,m}$ solve $G_n(t)=m/n$ and put
$z_{n,m}:=\log(n/m)$.  Uniformly for
$(\log n)^2\le m\le\delta_n n$, where
$\delta_n=(\log\log n)^{-4}$,
\begin{align}
 t^0_{n,m}&=\frac{m}{C_+n z_{n,m}}
 \left(1+O\left(\frac{\log z_{n,m}}{z_{n,m}}
 +(m/n)^\eta+m^{-1}\right)\right),
 \label{eq:critical-time}\\
 \frac{\theta_{n,m}}{nD_n(t^0_{n,m})}
 &=\frac{d}{m z_{n,m}}(1+o(1)),
 \qquad
 \theta_{n,m}t^0_{n,m}=\frac{d}{z_{n,m}}(1+o(1)).
 \label{eq:critical-hazard}
\end{align}
The fast-corner concentration proof applies with the regularly varying
factor $\log(1/t)$ retained.  Consequently, if
$B^0_{n,m}:=\one\{E_{n,m}\ge t^0_{n,m}\}$, then
\begin{equation}\label{eq:critical-compensator-approximation}
 \E\sum_{m=(\log n)^2}^{\lfloor\delta_n n\rfloor}
 \left|p_{n,m}-
 \frac{\theta_{n,m}}{nD_n(t^0_{n,m})}B^0_{n,m}\right|=o(1).
\end{equation}
The $B^0_{n,m}$ are independent, their means are
$1+O(1/z_{n,m})$, and the sum of the squared deterministic coefficients is
$o(1)$.  Hence
\begin{align}
 \sum_{m=(\log n)^2}^{\lfloor\delta_n n\rfloor}p_{n,m}
 &=d\sum_{m=(\log n)^2}^{\lfloor\delta_n n\rfloor}
 \frac1{m\log(n/m)}+o_{\Pp}(1)\\
 &=d\log\log n+o_{\Pp}(\sqrt{\log\log n}).
 \label{eq:critical-main-compensator}
\end{align}

It remains to remove the two cutoffs.  Before $(\log n)^2$ draws, at most
$(\log n)^2$ clocks have been deleted.  The untouched part of the left
power block has total rate at least $cn\log(n/(\log n)^2)$, whereas
$\theta_{n,m}\le Cn/m$.  Therefore
\[
 \sum_{m< (\log n)^2}p_{n,m}
 \le\frac{C}{\log(n/(\log n)^2)}
 \sum_{m<(\log n)^2}\frac1m=o(1)
\]
deterministically.  For $m>\delta_n n$, the corresponding one-point estimate
$\E p_{n,m}\le C/[m\log(n/m)]$ gives a total contribution
$O(\log\log(1/\delta_n))=o(\sqrt{\log\log n})$.  Thus the full fixed-point
compensator is
\[
 d\log\log n+o_{\Pp}\bigl(\sqrt{\log\log n}\bigr).
\]
Its square sum is $o_{\Pp}(\log\log n)$ by the same bounds.  Proposition
\ref{prop:conditional-clt}, with $v_n=d\log\log n$, proves
\eqref{eq:critical-fixed-clt}; taking expectations in
the preceding estimates proves the mean asymptotic.

For $\ell\ge2$, the critical logarithmic increment density has parameter
$q_m\asymp1/z_{n,m}$.  Write $p_q:=p_{1,q}$.  Direct convolution of the
representation $W=\log q_m-\log Y$ gives, for $0<q\le1/2$,
\[
 p_q^{*\ell}(0)\le C_\ell q^\ell
 (1+|\log q|)^{\ell-1}.
\]
The corresponding rooted trace is bounded by
\[
 \frac{C_\ell}{m}
 \frac{(1+\log z_{n,m})^{\ell-1}}{z_{n,m}^\ell},
\]
whose sum is finite for $\ell\ge2$.  The cutoff ranges are controlled as
above.  In the simple profile of Theorem~\ref{thm:critical-pole}, rates
outside the critical left window are bounded above and below; the usual
bounded-rate insertion estimate gives a uniformly bounded contribution
there as well.  Hence $\sup_n\E C_{n,\ell}<\infty$, completing the proof.
\end{proof}

\subsection{Cycles in the endpoint corners}

We now pass from one-point compensators to finite cycle closures.  First we
dispose of cycles rooted a macroscopic distance from either corner.

\begin{lemma}\label{lem:power-compact-cycles}
Under Assumptions~\ref{ass:power-endpoints}
and~\ref{ass:power-population}, for every $\delta>0$ and fixed $L$,
\begin{equation}\label{eq:power-compact-cycle-bound}
 \sup_n\E\sum_{\ell=1}^L
 \#\{\ell\text{-cycles whose largest vertex lies in }
 [\delta n,(1-\delta)n]\}<\infty.
\end{equation}
\end{lemma}

\begin{proof}
The endpoint expansions leave a linear reservoir of rates in some fixed
interval $[d,M]$ after any fixed number of deletions.  Consequently, for
every fixed rank interval $[\delta,1-\delta]$, the corresponding order
times stay in a deterministic compact subinterval of $(0,\infty)$ with
exponentially high probability.
Conditional on the background race, insertion into a middle gap is therefore
bounded by
\[
 \frac Cn\theta_{n,i}e^{-c\theta_{n,i}}\le\frac{C'}n.
\]
The product bound holds for any fixed collection of distinct marked clocks.
Summing it over the $O(n^\ell)$ cyclic assignments proves the result.  The
exceptional order-time event is absorbed by its exponentially small
probability.
\end{proof}

At an active right corner, uniformly for $u,v\downarrow0$ with bounded
ratio, the local rank kernel is
\begin{equation}\label{eq:corner-kernel}
 \rho(1-u,1-v)\sim\frac1v k_R(u/v),
 \qquad
 k_R(r):=\beta q_Rr^\beta e^{-q_Rr^\beta}.
\end{equation}
At an active left corner it is
\begin{equation}\label{eq:left-corner-kernel}
 \rho(u,v)\sim\frac1v k_L(u/v),
 \qquad
 k_L(r):=\alpha q_Lr^{-\alpha}e^{-q_Lr^{-\alpha}}.
\end{equation}
In logarithmic distance these are translation-invariant Markov kernels with
increment densities $p_{\beta,q_R}$ and the reflection of
$p_{\alpha,q_L}$.  We now use these kernels to follow all vertices of a cycle
of fixed length.

Let
\[
 S_{n,k-1}:=\{\pi_n(1),\ldots,\pi_n(k-1)\}.
\]
Starting at $v_0=k$, follow the exposed partial rank map: if
$v_a\in S_{n,k-1}$, put $v_{a+1}=R_{n,v_a}$.  Stop at the first
$d=d_{n,k}$ for which $v_d\notin S_{n,k-1}$, and set $d_{n,k}=\infty$ if a
vertex first repeats.  Define
\begin{equation}\label{eq:cycle-closure-probability}
 q_{n,k,\ell}:=
 \one\{d_{n,k}=\ell-1\}\frac{\theta_{n,v_{\ell-1}}}{W_{n,k}},
 \qquad
 J_{n,k,\ell}:=
 \one\{d_{n,k}=\ell-1,\ \pi_n(k)=v_{\ell-1}\}.
\end{equation}

\begin{lemma}\label{lem:cycle-closure}
For every $\ell\ge1$,
\begin{equation}\label{eq:cycle-closure-sum}
 C_{n,\ell}=\sum_{k=1}^nJ_{n,k,\ell},
 \qquad
 \E(J_{n,k,\ell}\mid\cF_{n,k-1})=q_{n,k,\ell}.
\end{equation}
For each $k$, at most one of the numbers $q_{n,k,\ell}$ is nonzero.
\end{lemma}

\begin{proof}
At draw $k$, selecting a remaining label $i$ adds the edge $i\mapsto k$ to
the partial rank map.  This closes a cycle exactly when the exposed path
starting at $k$ first reaches $i$.  If it does so after $\ell-1$ exposed
edges, the cycle has length $\ell$.  Conversely, every cycle is exposed
exactly once, when the edge whose target is its largest vertex is added.
This proves the first identity.  Given $\cF_{n,k-1}$, the endpoint and
remaining weight are known, and the Luce rule selects the endpoint with the
probability in \eqref{eq:cycle-closure-probability}.  The first undrawn
endpoint has a unique depth.
\end{proof}

For $\ell=1$, this is exactly the predictable fixed-point probability
$p_{n,k}$.  The closure identity also gives a convenient vector martingale
criterion.

\begin{proposition}\label{prop:vector-closure-clt}
Fix $L$.  Suppose that, for positive constants $b_1,\ldots,b_L$,
\begin{equation}\label{eq:vector-closure-input}
 \sum_{k=1}^nq_{n,k,\ell}
 =b_\ell\log n+o_{\Pp}(\sqrt{\log n}),
 \qquad
 \sum_{k=1}^nq_{n,k,\ell}^2=o_{\Pp}(\log n)
\end{equation}
for every $\ell\le L$.  Then
\[
 \left(
 \frac{C_{n,\ell}-b_\ell\log n}{\sqrt{b_\ell\log n}}
 \right)_{\ell=1}^L
 \Longrightarrow (Z_1,\ldots,Z_L),
\]
where the $Z_\ell$ are independent standard normal variables.
\end{proposition}

\begin{proof}
Put
$M_{n,\ell}=\sum_k(J_{n,k,\ell}-q_{n,k,\ell})$.
Its predictable variance is
$\sum_kq_{n,k,\ell}(1-q_{n,k,\ell})\sim b_\ell\log n$.
If $\ell\ne r$, Lemma~\ref{lem:cycle-closure} gives both
$J_{n,k,\ell}J_{n,k,r}=0$ and
$q_{n,k,\ell}q_{n,k,r}=0$, so the predictable cross variation vanishes.
The increments are bounded, and the multivariate martingale central limit
theorem~\cite{hall1980martingale} proves the claim.
\end{proof}

It remains to control the open-path compensators.  Put
\[
 Q^R_{n,\ell}(A,B):=\sum_{m=A}^Bq_{n,k_m,\ell},
 \qquad
 Q^L_{n,\ell}(A,B):=\sum_{m=A}^Bq_{n,m,\ell},
 \qquad T(A,B):=\log(B/A).
\]

\begin{lemma}\label{lem:power-two-path}
For every fixed $L$ there are $\kappa>0$, $\delta_0>0$, and $C<\infty$ such
that, for $\ell\le L$, $1\ll A\le B\le\delta_0n$, and every active corner
$\sigma\in\{L,R\}$,
\begin{align}
 \left|\E Q^\sigma_{n,\ell}(A,B)-b_{\sigma,\ell}T(A,B)\right|
 &\le C\left(A^{-\kappa}+(B/n)^\kappa(1+T(A,B))\right),
 \label{eq:path-first-moment}\\
 \Var Q^\sigma_{n,\ell}(A,B)
 &\le C\left(A^{-\kappa}(1+T(A,B))
 +(B/n)^\kappa(1+T(A,B))^2\right),
 \label{eq:path-second-moment}\\
 \E\sum_{m=A}^B(q^\sigma_{n,m,\ell})^2&\le CA^{-\kappa},
 \label{eq:path-square-bound}
\end{align}
where $q^R_{n,m,\ell}:=q_{n,k_m,\ell}$ and
$q^L_{n,m,\ell}:=q_{n,m,\ell}$.  If $\delta_n\downarrow0$ and
$\delta_n n\to\infty$, then
\begin{equation}\label{eq:path-outer}
 \E\sum_{\delta_n n<k<(1-\delta_n)n}q_{n,k,\ell}
 \le C(1+\log(1/\delta_n)).
\end{equation}
For the cutoffs used below, an inactive endpoint window has expected
compensator $O(\log\log n)$; the only borderline case is a left pole of
order one, which has this order already for fixed points.
\end{lemma}

\begin{proof}
We combine the two endpoint population estimates with the finite insertion
paths of Lemma~\ref{lem:finite-insertion-path}.  We give the right-corner
calculation first and then the small-time left-corner calculation.

\smallskip
\noindent\emph{Step 1: the exact gap formula.}
Delete at most $2L$ marked clocks and consider the resulting background
race.  Write $s_h^\circ$ for the left endpoint of a gap at which $h$
background clocks remain and $W_h^\circ$ for their total rate.  Conditional
on the elimination chain,
\begin{equation}\label{eq:tail-gap-exact}
 \Delta_h=\frac{\xi_h}{W_h^\circ},
 \qquad \xi_h\sim\operatorname{Exp}(1),
\end{equation}
and the variables $\xi_h$ belonging to distinct gaps are independent.  Thus
the exact conditional probability that an independent clock of rate
$\theta_{n,k_a}$ enters this gap is
\begin{equation}\label{eq:one-gap-insertion-exact}
 e^{-\theta_{n,k_a}s_h^\circ}
 \E\left(1-e^{-\theta_{n,k_a}\Delta_h}
          \mid\text{elimination chain}\right)
 =e^{-\theta_{n,k_a}s_h^\circ}
   \frac{\theta_{n,k_a}}{W_h^\circ+\theta_{n,k_a}}.
\end{equation}
This identity is the reason the population estimate must control both the
gap location and the remaining rate.

Assume first that the right corner is active and abbreviate
$\mu_h=nD_n(t_{n,h}^-)$.  Equations
\eqref{eq:random-time-concentration}--\eqref{eq:random-weight-at-order}, and
their unchanged versions after a bounded deletion, give an event $G_h$ on
which
\begin{equation}\label{eq:path-population-good}
 s_h^\circ=t_{n,h}^-(1+O(h^{-1/4})),\qquad
 W_h^\circ=\mu_h(1+O(h^{-1/4})),
 \qquad
 \Pp(G_h^c)\le Ch^{-4}+Ce^{-c\sqrt h}.
\end{equation}
For $a/h$ in a fixed compact subset of $(0,\infty)$,
Proposition~\ref{prop:deterministic-slow} and
\eqref{eq:one-gap-insertion-exact} therefore give
\begin{equation}\label{eq:discrete-tail-kernel}
 \overline K_n(a,h)
 =\frac1h k_R(a/h)
 \left(1+O(h^{-1/4}+(h/n)^\xi+h^{-1})\right),
\end{equation}
where $\overline K_n(a,h)$ denotes the insertion probability averaged over
the background.  The estimate holds jointly for any fixed collection of
gaps, because the union of their bad events has the same order.

\smallskip
\noindent\emph{Step 2: removal of the ratio cutoff.}
The approximation in \eqref{eq:discrete-tail-kernel} must be summed over
unbounded ratios.  The rough tail scales and $xe^{-x}\le Ce^{-x/2}$ for
$x\ge1$ give, uniformly in $a,h$ in the tail,
\begin{equation}\label{eq:tail-kernel-envelope}
 \overline K_n(a,h)\le\frac Ch\,\omega(a/h),\qquad
 \omega(r):=r^\beta\one\{r<1\}
       +r^\beta e^{-cr^\beta}\one\{r\ge1\}.
\end{equation}
The same bound holds with any bounded number of clocks deleted.  Notice that
$\omega(e^z)$ is integrable in $z$ and has an exponential moment of some
positive order in both tails.  We may therefore first apply
\eqref{eq:discrete-tail-kernel} on a bounded logarithmic window and then let
the window grow as a small multiple of $\log(a\wedge h)$.  After decreasing
$\kappa>0$ if necessary, this yields the summable error
\begin{equation}\label{eq:uniform-tail-kernel-error}
 \left|\overline K_n(a,h)-\frac1h k_R(a/h)\right|
 \le\frac Ch\left((a\wedge h)^{-\kappa}
        +((a\vee h)/n)^\kappa\right)\widetilde\omega(a/h),
\end{equation}
for an envelope $\widetilde\omega$ with the same integrability and slightly
smaller exponential moments.  If a proposed path leaves the right endpoint
window, its first macroscopic crossing costs $O((a/n)^\kappa)$ after
summation.  Splitting the marked rate at $1/s_h^\circ$ and using the global
population laws in Assumption~\ref{ass:power-population} gives this bound
without any pointwise interior assumption.  A path reaching a fast left
corner must also cross back to close near the right endpoint, and has the
same summable cost.

\smallskip
\noindent\emph{Step 3: one rooted path.}
Expand the open path in \eqref{eq:cycle-closure-probability}:
\begin{align}
 q_{n,k,\ell}
 =\sum_{\boldsymbol v}
 \frac{\theta_{n,v_{\ell-1}}}{W_{n,k}}\,
 &\one\{R_{n,k}=v_1,\ R_{n,v_1}=v_2,\ldots,
          R_{n,v_{\ell-2}}=v_{\ell-1}\}
 \one\{R_{n,v_{\ell-1}}\ge k\},
 \label{eq:open-path-expansion}
\end{align}
where the sum is over distinct vertices below $k$; for $\ell=1$ the internal
path is empty.  Swap the distinct replacement clocks of one proposed
path into the background.  The swap preserves their joint law and, as in
Lemma~\ref{lem:finite-insertion-path}, turns the proposal into a genuine path
in a background permutation, up to a bounded displacement of its prescribed
ranks.  Equations \eqref{eq:uniform-tail-kernel-error} and
\eqref{eq:tail-kernel-envelope}, applied successively along this fixed-depth
path, give
\begin{align}
 \E q_{n,k_m,\ell}
 &=\sum_{\substack{m_1,\ldots,m_{\ell-1}>m\\\mathrm{distinct}}}
 \prod_{j=0}^{\ell-1}
 \frac1{m_{j+1}}k_R(m_j/m_{j+1})
 +O\left(m^{-1-\kappa}+m^{-1}(m/n)^\kappa\right),
 \label{eq:discrete-rooted-trace}
\end{align}
where $m_0=m_\ell=m$.  The bounded displacement and the exclusion of
coincident internal vertices are absorbed by the first error term.

Put $z_j=\log(m_j/m)$.  Riemann summation in
\eqref{eq:discrete-rooted-trace}, justified by
\eqref{eq:tail-kernel-envelope}, gives
\begin{equation}\label{eq:rooted-trace-integral}
 \E q_{n,k_m,\ell}
 =\frac1m\int_{(0,\infty)^{\ell-1}}
 \prod_{j=0}^{\ell-1}p_{\beta,q_R}(z_j-z_{j+1})
 \,dz_1\cdots dz_{\ell-1}
 +O\left(m^{-1-\kappa}+m^{-1}(m/n)^\kappa\right),
\end{equation}
with $z_0=z_\ell=0$.  For i.i.d. continuous increments conditioned to sum
to zero, exactly one of the $\ell$ cyclic shifts has its initial point as the
unique minimum.  The cyclic lemma therefore identifies the integral in
\eqref{eq:rooted-trace-integral} as
$p_{\beta,q_R}^{*\ell}(0)/\ell=b_{R,\ell}$.  We have proved the pointwise
estimate
\begin{equation}\label{eq:one-path-pointwise}
 \E q_{n,k_m,\ell}
 =\frac{b_{R,\ell}}m
 +O\left(m^{-1-\kappa}+m^{-1}(m/n)^\kappa\right).
\end{equation}

\smallskip
\noindent\emph{Step 4: two paths.}
Use two independent families of replacement clocks.  If the two proposed
paths have disjoint vertex sets, swap both families into the common
background.  Conditional on its elimination chain, distinct gaps use
independent variables in \eqref{eq:tail-gap-exact}; hence the leading
one-gap terms factor.  The remaining covariance comes from the population
errors in \eqref{eq:uniform-tail-kernel-error}.  If the paths intersect,
fix their first common vertex.  The bounded approximate indegree in
Lemma~\ref{lem:finite-insertion-path} leaves at most a bounded number of
predecessors and supplies one additional factor $(m\wedge m')^{-1}$.
Summing over the two fixed depths gives
\begin{align}
 \left|\operatorname{Cov}(q_{n,k_m,\ell},q_{n,k_{m'},\ell})\right|
 &\le \frac{C}{mm'}\left((m\wedge m')^{-\kappa}
       +((m\vee m')/n)^\kappa\right)
   +\frac{C}{m^2}\one\{m=m'\},\label{eq:two-path-pointwise}\\
 \E q_{n,k_m,\ell}^2
 &\le Cm^{-2}+Cm^{-4}+Ce^{-c\sqrt m}.
 \label{eq:path-square-pointwise}
\end{align}
The last two terms in the second line absorb the complement of
\eqref{eq:path-population-good}.  This is the only point where two-path
collisions enter; the distinct-source structure of a genuine cycle is what
allows the simultaneous swap.

\smallskip
\noindent\emph{Step 5: the left corner.}
Assume the left corner is active.  At the gap for rank $h$, Propositions
\ref{prop:deterministic-early} and~\ref{prop:early-compensator} turn
\eqref{eq:one-gap-insertion-exact} into
\begin{equation}\label{eq:discrete-left-kernel}
 \overline K_n^L(a,h)=\frac1h k_L(a/h)
 \left(1+O(h^{-1/4}+(h/n)^\xi+h^{-1})\right).
\end{equation}
The logarithmic envelope is the reflection of
\eqref{eq:tail-kernel-envelope}.  A cycle closed at $m$ has all its other
vertices below $m$, so its open path cannot leave the left endpoint window.
Repeating the one- and two-path swaps with $m_j<m$ gives
\begin{align}
 \E q_{n,m,\ell}
 &=\frac{b_{L,\ell}}m
 +O\left(m^{-1-\kappa}+m^{-1}(m/n)^\kappa\right),
 \label{eq:left-one-path-pointwise}\\
 \left|\operatorname{Cov}(q_{n,m,\ell},q_{n,m',\ell})\right|
 &\le \frac{C}{mm'}\left((m\wedge m')^{-\kappa}
 +((m\vee m')/n)^\kappa\right)
 +\frac C{m^2}\one\{m=m'\},
 \label{eq:left-two-path-pointwise}\\
 \E q_{n,m,\ell}^2
 &\le Cm^{-2}+Cm^{-4}+Ce^{-c\sqrt m}.
 \label{eq:left-path-square-pointwise}
\end{align}
Indeed, with $z_j=\log(m_j/m)<0$, the cyclic lemma now selects the unique
maximum.  Reflection does not change an $\ell$-fold return density at zero,
so the trace is $p_{\alpha,q_L}^{*\ell}(0)/\ell=b_{L,\ell}$.

Summing \eqref{eq:one-path-pointwise} gives
\eqref{eq:path-first-moment} for the right corner, and summing
\eqref{eq:left-one-path-pointwise} gives it for the left corner.  Likewise,
summing \eqref{eq:two-path-pointwise} and
\eqref{eq:left-two-path-pointwise}, and using
\[
\begin{aligned}
 \sum_{A\le m,m'\le B}\frac{(m\wedge m')^{-\kappa}}{mm'}
 &\le CA^{-\kappa}(1+T(A,B)),\\
 \sum_{A\le m,m'\le B}
 \frac{((m\vee m')/n)^\kappa}{mm'}
 &\le C(B/n)^\kappa(1+T(A,B))^2,
\end{aligned}
\]
proves \eqref{eq:path-second-moment}.  Finally,
\eqref{eq:path-square-pointwise} and
\eqref{eq:left-path-square-pointwise} sum to $O(A^{-1})$; decreasing
$\kappa$ to at most one proves \eqref{eq:path-square-bound}.

For \eqref{eq:path-outer}, the one-path envelopes give $C/k$ and
$C/(n-k)$ in the two transition bands, while
Lemma~\ref{lem:power-compact-cycles} handles a fixed middle interval.
At a subdominant zero or pole, the exponential argument tends to infinity
or zero by a fixed power, so its endpoint sum is bounded.  A pole of order
one instead gives $O(1/(m\log(n/m)))$ for fixed points and a summable bound
for longer cycles.  This proves both the middle estimate and the stated
$O(\log\log n)$ inactive-window bound.
\end{proof}

\begin{proof}[Proof of Theorem~\ref{thm:power-spatial}]
Write $I_j=(a_j,b_j]$.  By the closure identity,
\begin{align*}
 C_{n,\ell}^{L}(I_j)
 &=\sum_{n^{a_j}<m\le n^{b_j}}J_{n,m,\ell},&
 C_{n,\ell}^{R}(I_j)
 &=\sum_{n^{a_j}<m\le n^{b_j}}J_{n,k_m,\ell}.
\end{align*}
The corresponding predictable compensators are the same sums with $J$
replaced by $q$.  Apply Lemma~\ref{lem:power-two-path} with
$A=\lfloor n^{a_j}\rfloor$ and $B=\lfloor n^{b_j}\rfloor$.  Because
$0<a_j<b_j<1$, its error terms tend to zero, and hence, for every active
corner $\sigma$,
\begin{align}
 \sum_{n^{a_j}<m\le n^{b_j}}q^\sigma_{n,m,\ell}
 &=b_{\sigma,\ell}(b_j-a_j)\log n+o_{\Pp}(1),
 \label{eq:spatial-compensator}\\
 \sum_{n^{a_j}<m\le n^{b_j}}
 (q^\sigma_{n,m,\ell})^2&=o_{\Pp}(1).
 \label{eq:spatial-square-sum}
\end{align}
The first estimate also holds in mean.

Center each sum of $J$ by its predictable compensator.  These are bounded
martingales indexed by the sequential exposure.  Two coordinates belonging
to disjoint logarithmic intervals or different corners use disjoint draw
indices.  At a common draw index, coordinates belonging to different cycle
lengths have zero product by Lemma~\ref{lem:cycle-closure}.  Thus all
off-diagonal predictable covariations vanish, while
\eqref{eq:spatial-compensator}--\eqref{eq:spatial-square-sum} give the stated
diagonal variances.  The multivariate martingale central limit theorem proves
\eqref{eq:power-spatial-clt}.  Taking expectations proves
\eqref{eq:power-spatial-mean}.  The localization assertion is immediate for
$\ell=1$.  For
$\ell\ge2$, the logarithmic envelopes in
\eqref{eq:tail-kernel-envelope} and its left-corner reflection have an
exponential moment.  Requiring one of the at most $\ell-1$ path increments
to exceed $\delta\log n/(\ell-1)$ therefore adds a factor $O(n^{-c\delta})$
to the one-path sum, uniformly for closure distances in $n^I$.
Summing over that compact logarithmic annulus proves
the localization assertion following Theorem~\ref{thm:power-spatial}.
\end{proof}

\begin{proof}[Completion of the proof of Theorem~\ref{thm:power-clt}]
Let $\kappa$ be supplied by Lemma~\ref{lem:power-two-path}, and take
\[
 A_n=(\log n)^{1/4},\qquad
 \delta_n=(\log n)^{-K},
\]
where $K\kappa>2$.  Equations \eqref{eq:path-first-moment} and
\eqref{eq:path-second-moment}, followed by Chebyshev's inequality, give
\begin{equation}\label{eq:cycle-compensator-corners}
 Q^\sigma_{n,\ell}(A_n,\lfloor\delta_n n\rfloor)
 =b_{\sigma,\ell}\log(\delta_n n/A_n)
 +o_{\Pp}(\sqrt{\log n})
\end{equation}
for every active corner $\sigma$ and fixed $\ell$.  The two ranges with
$m<A_n$ contribute at most $2A_n=o(\sqrt{\log n})$.  The middle range and
all inactive endpoint contributions are $o_{\Pp}(\sqrt{\log n})$ by
Lemma~\ref{lem:power-two-path} and Markov's inequality.  Since
$\log(\delta_n/A_n)=O(\log\log n)$, summing the active corners gives
\[
 \sum_{k=1}^nq_{n,k,\ell}
 =B_\ell\log n+o_{\Pp}(\sqrt{\log n}).
\]

On the endpoint ranges, \eqref{eq:path-square-bound} and Markov's inequality
give a square sum $o_{\Pp}(\log n)$.  On the final and middle ranges,
use $q_{n,k,\ell}^2\le q_{n,k,\ell}$ together with the preceding bounds.
Thus \eqref{eq:vector-closure-input} holds with
$b_\ell=B_\ell$, and Proposition~\ref{prop:vector-closure-clt}
proves the joint central limit theorem.

Finally, take expectations in \eqref{eq:path-first-moment}.  The final,
middle, and inactive ranges are $o(\log n)$, so
$\E C_{n,\ell}=\E\sum_kq_{n,k,\ell}\sim B_\ell\log n$.
\end{proof}

\section{Independent-clock extensions}\label{sec:clock-extensions}

The bulk density~\eqref{eq:rho} uses the exponential form of the Luce race,
but the last-position estimate has a distribution-free core.  This identifies
a broader class of ranking models for which the same fixed-point strategy is
available.

Let $T_1,\ldots,T_n$ be independent continuous random variables with
densities $g_i$ and survival functions $\overline G_i$.  Order them
increasingly and write
\[
 R_i=1+\sum_{j\ne i}\one\{T_j<T_i\},\qquad
 S(t):=\sum_{i=1}^n\overline G_i(t).
\]
As before, put $k_m=n-m+1$.

\begin{proposition}
\label{prop:general-clock-tail}
Let $1\le M\le n$, choose $B$ with $B-1\ge2M$, and let $s$ satisfy
$S(s)=B$.  Suppose that an integrable function $h$ satisfies
\begin{equation}\label{eq:clock-density-envelope}
 g_{k_m}(t)\le h(t),\qquad 1\le m\le M,\quad t\ge s.
\end{equation}
Then, for a universal $c>0$,
\begin{equation}\label{eq:general-clock-tail}
 \sum_{m=1}^M\Pp(R_{k_m}=k_m)
 \le M e^{-cB}+2\int_s^\infty h(t)\,dt.
\end{equation}
\end{proposition}

\begin{proof}
Conditioning on $T_{k_m}=t$ gives the exact identity
\[
 \Pp(R_{k_m}=k_m)=\int_{-\infty}^{\infty}g_{k_m}(t)
 \Pp\left(\sum_{j\ne k_m}\one\{T_j>t\}=m-1\right)dt.
\]
For $t\le s$, the survivor count in parentheses has mean at least
$S(t)-1\ge B-1\ge2M$, so its lower tail is at most $e^{-cB}$ by a Chernoff
bound.  The integral over this range, summed over $m$, is at most
$Me^{-cB}$.

For fixed $t$, couple all the inner probabilities using one independent
sample of the clocks.  If $N(t)=\sum_j\one\{T_j>t\}$, the event in the last
display can hold only for $m=N(t)$ or $m=N(t)+1$.  Hence the sum of its
indicators over $m\le M$ is at most two.  On $t\ge s$, use
\eqref{eq:clock-density-envelope}, take expectations, and integrate to obtain
the second term in~\eqref{eq:general-clock-tail}.
\end{proof}

Suppose more generally that the clock attached to label $i/n$ has a limiting
distribution function $G(x,t)$ and density $g(x,t)$.  If
\[
 F(t):=\int_0^1G(x,t)\,dx,\qquad
 d(t):=\int_0^1g(x,t)\,dx,
\]
assume that $F$ is strictly increasing on the range under consideration and
that $d(t_y)>0$.  With $t_y=F^{-1}(y)$, the natural bulk density is
\begin{equation}\label{eq:general-clock-density}
 \rho_T(x,y)=\frac{g(x,t_y)}{d(t_y)}.
\end{equation}
Thus a bulk local-rank theorem together with
Proposition~\ref{prop:general-clock-tail} gives the same architecture as the
fixed-point proof above, provided one can choose triangular-array parameters
$M_n,B_n,s_n,h_n$ for which the right side of
\eqref{eq:general-clock-tail} tends to zero.  A candidate bulk density is not
enough by itself: one must still prove compensator convergence, or an
equivalent microscopic local law.

Longer cycles require another input.  Their late closing label need not be a
terminal label, so Proposition~\ref{prop:general-clock-tail} alone does not
control them.  In the exponential case, Lemma~\ref{lem:finite-insertion-path}
and Propositions~\ref{prop:cycle-rate-truncation}
and~\ref{prop:cycle-endpoint-bound} overcome this obstruction by swapping a
finite insertion path into the background, charging exceptional rates, and
marking the edge into the largest cycle label.  The finite-path lemma is
distribution-free; the truncation and shell bounds for another clock family
must come from suitable density comparisons and endpoint envelopes.

A concrete testing ground is the common-shape Gamma ranking family of
\cite{stern1990models}.  Shape one is the exponential race and hence the Luce
model.  For fixed Gamma shape and terminal rates bounded below, the terminal
densities have a common polynomial-times-exponential envelope, so
Proposition~\ref{prop:general-clock-tail} applies; a bulk local law remains to
be established.  The independent-utility ranking models of
\cite{thurstone1927law} are another natural target, provided their terminal
densities admit an integrable envelope and the corresponding bulk local limit
is available.  Separately, the finite-dimensional criterion of
\cite{mukherjee2016fixed} already yields the cyclic trace law
\eqref{eq:cycle-intensity} for permuton-sampled permutations with strictly
positive continuous densities, Mallows permutations with $n(1-q_n)$
convergent, and certain continuous exponential-family models.

The extension also has sharp limits.  For each fixed Ewens
parameter $\vartheta>0$, asymptotically uniform one-point marginals coexist
with limiting $\ell$-cycle mean $\vartheta/\ell$~\cite{tavare2021magical};
thus a microscopic hypothesis remains essential.  For fixed $q\in(0,1)$,
Mallows permutations lie in a localized regime in which bounded-cycle counts
have linear means and joint Gaussian rather than Poisson
fluctuations~\cite{he2023cycles}.  These examples underline the central
point: the macroscopic limit predicts the form of the cycle constants only
after a model-specific argument has reached the microscopic scale.

\bibliographystyle{alphaabbr}
\small
\bibliography{fixed_points}

\end{document}
